% \iffalse
%  Local Variables:
%  mode: doctex
%  TeX-master: t
%  End:
% \fi
%
% \iffalse meta-comment
%
% Copyright (C) 2010 by Hu,Weiyi <huweiyi@whut.edu.cn>
%
% This file may be distributed and/or modified under the
% conditions of the LaTeX Project Public License, either version 1.3a
% of this license or (at your option) any later version.
% The latest version of this license is in:
%
% http://www.latex-project.org/lppl.txt
%
% and version 1.3a or later is part of all distributions of LaTeX
% version 2004/10/01 or later.
%
% Created at 2010-03-02 15:21
% This template is based on the work of thuthesis by Ruini Xue,
% \fi
%
% \CheckSum{0}
% \CharacterTable
%  {Upper-case    \A\B\C\D\E\F\G\H\I\J\K\L\M\N\O\P\Q\R\S\T\U\V\W\X\Y\Z
%   Lower-case    \a\b\c\d\e\f\g\h\i\j\k\l\m\n\o\p\q\r\s\t\u\v\w\x\y\z
%   Digits        \0\1\2\3\4\5\6\7\8\9
%   Exclamation   \!     Double quote  \"     Hash (number) \#
%   Dollar        \$     Percent       \%     Ampersand     \&
%   Acute accent  \'     Left paren    \(     Right paren   \)
%   Asterisk      \*     Plus          \+     Comma         \,
%   Minus         \-     Point         \.     Solidus       \/
%   Colon         \:     Semicolon     \;     Less than     \<
%   Equals        \=     Greater than  \>     Question mark \?
%   Commercial at \@     Left bracket  \[     Backslash     \\
%   Right bracket \]     Circumflex    \^     Underscore    \_
%   Grave accent  \`     Left brace    \{     Vertical bar  \|
%   Right brace   \}     Tilde         \~}
%
% \changes{v0.1.0}{2010/03/02}{初始版本}
% \iffalse
%<*driver>
\ProvidesFile{WHUTthesis.dtx}[2010/06/19 0.1.2d Thesis Template for Wuhan University of Technology]
\documentclass[10pt]{ltxdoc}
\usepackage{dtx-style}
\EnableCrossrefs
\CodelineIndex
\RecordChanges
%\OnlyDescription
\begin{document}
    \DocInput{\jobname.dtx}
\end{document}
%</driver>
% \fi
%
% \GetFileInfo{\jobname.dtx}
% \MakeShortVerb{\|}
%
% \def\whutthesis{\textsc{Whut}\-\textsc{Thesis}}
% \def\pkg#1{\texttt{#1}}
%

%
% \DoNotIndex{\begin,\end,\begingroup,\endgroup}
% \DoNotIndex{\ifx,\ifdim,\ifnum,\ifcase,\else,\or,\fi}
% \DoNotIndex{\let,\def,\xdef,\newcommand,\renewcommand}
% \DoNotIndex{\expandafter,\csname,\endcsname,\relax,\protect}
% \DoNotIndex{\Huge,\huge,\LARGE,\Large,\large,\normalsize}
% \DoNotIndex{\small,\footnotesize,\scriptsize,\tiny}
% \DoNotIndex{\normalfont,\bfseries,\slshape,\interlinepenalty}
% \DoNotIndex{\hfil,\par,\hskip,\vskip,\vspace,\quad}
% \DoNotIndex{\centering,\raggedright}
% \DoNotIndex{\c@secnumdepth,\@startsection,\@setfontsize}
% \DoNotIndex{\ ,\@plus,\@minus,\p@,\z@,\@m,\@M,\@ne,\m@ne}
% \DoNotIndex{\@@par,\DeclareOperation,\RequirePackage,\LoadClass}
% \DoNotIndex{\AtBeginDocument,\AtEndDocument}
%
% \IndexPrologue{\section*{索引}%
%    \addcontentsline{toc}{section}{索~~~~引}}
% \GlossaryPrologue{\section*{修改记录}%
%    \addcontentsline{toc}{section}{修改记录}}
%
% \renewcommand{\abstractname}{摘~~要}
% \renewcommand{\contentsname}{目~~录}
%
%
% \title{\whutthesis：武汉理工大学学位论文模板用户手册\thanks{Wuhan University of Technology \LaTeX{} Thesis Template.}}
% \author{{\fs 胡卫谊}\\[5pt]{\fs 武汉理工大学计算机科学与技术学院}\\[5pt] \texttt{huweiyi@whut.edu.cn}}
% \date{v\fileversion\ (\filedate)}
% \maketitle\thispagestyle{empty}
%
%
% \begin{abstract}\noindent
%   此宏包旨在建立一个简单易用的武汉理工大学学位论文模板，计划包括本科、硕士、
%   博士论文。现在已经支持理工类硕士和博士论文格式，对其它
%   格式的支持在后续维护中会陆续加入。
% \end{abstract}
%
% \vskip2cm
% \def\abstractname{免责声明}
% \begin{abstract}
% \noindent
% \begin{enumerate}
% \item 本模板的发布遵守 \LaTeX{} Project Public License，使用前请认真阅读协议内容。
% \item 本模板为作者根据武汉理工大学研究生院（筹）颁发的《博士、硕士学位论文书写格式的统一要求》
%   编写而成，旨在供武汉理工大学毕业生撰写学位
%   论文使用。
% \item 武汉理工大学研究生院（筹）只提供毕业论文写作指南，不提供官方模板，也不会授
%   权第三方模板为官方模板，所以此模板仅为写作指南的参考实现，不保证格式审查老师
%   不提意见。任何由于使用本模板而引起的论文格式审查问题均与本模板作者无关。
% \item 任何个人或组织以本模板为基础进行修改、扩展而生成的新的专用模板，请严格遵
%   守 \LaTeX{} Project Public License 协议。由于违犯协议而引起的任何纠纷争端均与
%   本模板作者无关。
% \end{enumerate}
% \end{abstract}
%
%
% \clearpage
% \begin{multicols}{2}[
%   \section*{\contentsname}
%   \setlength{\columnseprule}{.4pt}
%   \setlength{\columnsep}{18pt}]
%   \tableofcontents
% \end{multicols}
%
% \clearpage
% \pagenumbering{arabic}
% \pagestyle{headings}
% \section{模板介绍}
% \whutthesis 是为了帮助武汉理工大学毕业
% 生撰写毕业论文而编写的 \LaTeX{} 论文模板。
%
% 本文档将尽量完整地介绍模板的使用方法，如有不清楚之处可以参考示例文档或者给维护
% 人员（见后）写信，欢迎感兴趣的同学出力完善此使用手册。由于个人水平有限，虽然现
% 在的这个版本基本上满足了学校的要求，但难免还存在不足之处，欢迎大家积极反馈。
%
% {\color{blue}\fangsong 模板的作用在于减轻论文写作过程中格式调整的时间，其前提就是遵
%   守模板的用法，否则即使使用了 \whutthesis{} 也难以保证输出的论文符合学校规范。}
%
% {\color{red}\fangsong\heiti 注意：本模板在编写过程中在某些细节上偏向\LaTeX{}的排版美化约定，
% 因而与学校官方模板存在一些细微差异，请同学们参看给出的示例文档酌情使用！}
%
%
% \section{安装}
% \label{sec:installation}
%
% \subsection{下载}
% \whutthesis{} 主页：\url{http://whutthesis.googlecode.com/}。除此之外，不
% 再维护任何镜像。
%
% \whutthesis{} 的开发版本同样可以在 googlecode 上获得：
% \begin{shell}
% $ svn checkout http://whutthesis.googlecode.com/svn/trunk/ whutthesis-read-only 
% \end{shell}
%
% \subsection{模板的组成部分}
% 下图列出了模板的目录结构
%\begin{lstlisting}[title={模板目录结构},label=layout,numbers=none]
%  |-- WHUTthesis.ins
%  |-- WHUTthesis.dtx
%  |-- GBT7714-2005NLang.bst
%  |-- WHUTtils.sty
%  |-- msmake.cmd
%  |-- WHUTthesis.cls
%  |-- dtx-style.sty
%  |-- WHUTthesis.cfg
%  |-- mythesis.tex
%  |-- a3cover
%  |   |-- a3cover.tex
%  |   \-- spine.tex
%  |-- body
%  |   |-- abstract.tex
%  |   |-- ack.tex
%  |   |-- appendix01.tex
%  |   |-- chap01.tex
%  |   |-- chap02.tex
%  |   |-- cover.tex
%  |   |-- denotation.tex
%  |   \-- resume.tex
%  |-- figure
%  |   |-- whut.png
%  |   |-- whut.pdf
%  |   \-- ...
%  \-- reference
%      \-- refs.bib
%\end{lstlisting}
% 下表列出了 \whutthesis{} 的主要文件及其功能介绍：
%
% \begin{center}
%   \begin{longtable}{l|p{10cm}}
% \hline
% {\heiti 文件（夹）} & {\heiti 功能描述}\\\hline\hline
% \endfirsthead
% \hline
% {\heiti 文件（夹）} & {\heiti 功能描述}\\\hline\hline
% \endhead
% \endfoot
% \endlastfoot
% WHUTthesis.ins & 模板驱动文件 \\
% WHUTthesis.dtx & 模板文档代码的混合文件\\
% WHUTthesis.cls & 模板类文件\\
% WHUTthesis.cfg & 模板配置文件\\
% dtx-sty.sty & 用户手册样式文件\\
% GBT7714-2005NLang.bst & 参考文献样式文件\\\hline
% mythesis.tex & 示例文档主文件\\
% whuttils.sty & 为示例文档加载其它宏包\\
% a3cover/ & a3封面目录\\
% body/ & 示例文档章节具体内容目录\\
% figure/ & 示例文档图片路径目录\\
% reference/ & 示例文档参考文献目录\\\hline
% msmake.cmd & Windows 批处理工具\\\hline
% Readme & self-explanation\\
% \textbf{whutthesis.pdf} & 用户手册（本文档）\\\hline
%   \end{longtable}
% \end{center}
%
% 需要说明几点：1) WHUTthesis.cls、WHUTthesis.cfg、dtx-sty.sty和mythesis.tex可以
% 由 WHUTthesis.ins 和WHUTthesis.dtx 生成，但为了降低新手用户的使用难度，故将 cls
% 和 cfg 也一起发布。2) 学习一个新东西最好的办法就是读它的文
% 档：\emph{WHUTthesis.pdf}.
%
%
% \subsection{推荐的\TeX{}系统}
% \label{sec:recommend}
% 本模板当前版本v\fileversion\ (\filedate)在TexLive 2009下编写，未在其他\TeX{}系统上测试，
% 因此推荐用户使用TeXLive 2009。其安装包可以在\url{http://tug.org/texlive/}处下载。
% {\heiti 注意：}下载下来的iso压缩包可能需要支持xz解压缩的{\heiti7-Zip 4.65 稳定版 (2009-02-03)}
% 或更高版本。该版本可以从\url{http://7z.sparanoid.com/}处下载
%
% {\color{red}\zihao{4} 注意：由于本模板采用\XeLaTeX{}引擎处理，所以源文件编码应该
% 使用{\bf UTF-8}编码！}
%
% \subsection{开始安装}
% \label{sec:install}
%
% \subsubsection{生成模板}
% \label{sec:generate-cls}
% {\heiti 说明：默认的发行包中已经包含了所有文件，可以直接使用。如果对如何由 dtx 生
%   成模板文件以及模板文档不感兴趣，请跳过本小节。}
%
% 模板解压缩后生成文件夹 WHUTthesis-VERSION\footnote{VERSION 为版本号。}，其中包括：
% 模板源文件（WHUTthesis.ins 和 WHUTthesis.dtx），参考文献样式 GBT7714-2005NLang.bst，示例文档
% （mythesis.tex，WHUTtils.sty\footnote{我把可能用到但不一定用到的包以及一
%   些命令定义都放在这里面，以免 WHUTthesis.cls 过分臃
%   肿，减小增减依赖包时破坏模板的风险。}，body/ 和 figure/ 和 reference/）。在使用之前需要先生成模板文件和配置文件
% （具体命令细节请参考 |Readme| ）：
%
% \begin{shell}
% ~> cd WHUTthesis-VERSION
% # 仅生成模板（WHUTthesis.cls 和 WHUTthesis.cfg）
% ~/WHUTthesis-VERSION> msmake pkg
%
% # 如果模板和手册都需要生成，则
% ~/WHUTthesis-VERSION> msmake setup
%
% # 如果仅重新生成手册，则
% ~/WHUTthesis-VERSION> msmake man
% \end{shell}
%
% \subsubsection{生成自己的论文}
% \label{sec:my}
% 在编写自己的论文的时候，虽然可以通过GUI编辑器调用\XeLaTeX{}引擎来生成论文，但是频繁的点击按钮或者
% 输入命令仍然显得十分繁琐，所以\whutthesis 还提供了自动生成自己的论文的批处理工具\footnote{Linux平台下的
% makefile将在本模板的后续版本中提供}。
%
% \begin{shell}
% your_path $ msmake setup   # 生成宏包文件和说明文档
% your_path $ msmake my      # 生成示例文档
% your_path $ msmake a3      # 生成a3封面（在编译的时候可能出现错误，输入s跳过去就可以了）
% your_path $ msmake clean   # 清除临时文件
% \end{shell}
%
% \texttt{Makefile} 和 \texttt{msmake.cmd} 默认采用 \LaTeX\ 编译，可以根据自己的
% 需要修改命令。
%
%
% \subsection{升级}
% \label{sec:updgrade}
% \whutthesis{} 升级非常简单，下载最新的版本，
% 将 WHUTthesis.ins，WHUTthesis.dtx 和GBT7714-2005NLang.bst 拷贝至工作目录覆盖相应的文件，然后
% 运行：
% \begin{shell}
% $ latex WHUTthesis.ins
% \end{shell}
% 或者：
% \begin{shell}
% ~> msmake pkg
% \end{shell}
%
% 生成新的类文件和配置文件即可。当然也可以直接拷贝 WHUTthesis.cls, WHUTthesis.cfg
% 和 GBT7714-2005NLang.bst，免去上面命令的执行。只要明白它的工作原理，这个不难操作。
%
%
% \section{使用说明}
% \label{sec:usage}
% 本手册假定用户已经能处理一般的 \LaTeX{} 文档，并对 \BibTeX{} 有一定了解。如果你
% 从来没有接触过 \TeX 和 \LaTeX，建议先学习相关的基础知识。磨刀不误砍柴工！
%
% \subsection{关于提问}
% \label{sec:howtoask}
% 提问之前先问自己几个问题：
% \begin{enumerate}\addtolength{\itemsep}{-5pt}
% \item 我是不是认真地学习了 \LaTeX{} 基础知识？
% \item 我是不是认真地阅读了相关的文档？
% \item 我是不是 Google 了？
% \end{enumerate}
%
% 如果你确保自己已经完成了上面的操作，那么就可以到以下两个地方提问：
% \begin{itemize}\addtolength{\itemsep}{-5pt}
% \item 访问我们的wiki\url{http://whutthesis.googlecode.com} 
% \item 或直接给\href{mailto:whutthesis@gmail.com}{模板维护者}whutthesis@gmail.com写信。 
% \end{itemize}
%
%
% \subsection{\whutthesis{} 示例文件}
% \label{sec:userguide1}
% 模板核心文件只有三个：WHUTthesis.cls，WHUTthesis.cfg 和 GBT7714-2005NLang.bst，但是如果没有
% 示例文档用户会发现很难下手。所以推荐新用户从模板自带的示例文档入手，里面包括了
% 论文写作用到的所有命令及其使用方法，只需要用自己的内容进行相应替换就可以。对于
% 不清楚的命令可以查阅本手册。下面的例子描述了模板中章节的组织形式，来自于示例文
% 档，具体内容可以参考模板附带的 mythesis.tex 和 body/ 。
%
% \begin{example}
%    \begin{macrocode}
%<*myth>
%%%%%%%%%%%%%%%%%%%%%%%%%%%%%%%%%%%%%%%%%%%%%%%%%%%%%%%%%%%%%%%%%%%%%%
%%                                                                  %%
%%      武汉理工大学硕、博学位论文LaTeX模板主控文件(v0.1.0)         %%
%%      本模板建议在TeXLive2009下使用XeLaTex编译                    %%
%%                                                                  %%
%%                                                                  %%
%%                      胡卫谊 huweiyi@whut.edu.cn                  %%
%%                              2010/03/06                          %%
%%%%%%%%%%%%%%%%%%%%%%%%%%%%%%%%%%%%%%%%%%%%%%%%%%%%%%%%%%%%%%%%%%%%%%
%%
%% 1. 规范硕士导言
%%  \documentclass[master,oneside]{WHUTthesis}
%% 2. 规范博士导言（博士论文一般要双面打印）
%%  \documentclass[doctor,twoside,openright]{WHUTthesis}
%% 3. Adobe提供了免费的OTF字体，其页面显示效果较winfonts要好，但其只有宋体、
%%   仿宋、黑体、楷体等字体，缺少隶书和幼圆。如果启用OTF字体，则导言为：
%%  \documentclass[master,oneside,adobefonts]{WHUTthesis}
%%  \documentclass[doctor,twoside,openright,adobefonts]{WHUTthesis}
%%
\documentclass[master,oneside]{WHUTthesis}
%%
%% 所有其它可能用到的包都统一放到这里了，可以根据自己的实际需要添加或者删除。
\usepackage{WHUTtils}
%%
%% 你可以在这里修改配置文件中的定义，导言区可以使用中文。
%% \def\myname{胡卫谊}
%%
\begin{document}
%%
%% 定义图片文件目录
\graphicspath{{figure/}}
%%
%%
%% 封面（包括封面、独创性声明、授权书）
\frontmatter
%% 封面
%%%%%%%%%%%%%%%%%%%%%%%%%%%%%%%%%%%%%%%%%%%%%%%%%%%%%%%%%%%%%%%%%%%%%
%																	%
%		武汉理工大学硕、博学位论文LaTeX模板封面文件(v0.1.0)			%
%																	%
%																	%
%						胡卫谊 huweiyi@whut.edu.cn					%
%								2010/03/06							%
%%%%%%%%%%%%%%%%%%%%%%%%%%%%%%%%%%%%%%%%%%%%%%%%%%%%%%%%%%%%%%%%%%%%%
%%% Local Variables:
%%% mode: latex
%%% TeX-master: t
%%% End:
\categorynumber{123123}
%% 注明论文密级为公开、内部、秘密或机密
\secretlevel{公开}
\UDC{321321}

\ctitle{武汉理工大学学位论文 \LaTeX\ 模板使用示例文档}
\ctitleformated{武汉理工大学学位论文\LaTeX\ 模板使用示例文档}
% 根据自己的情况选，不用这样复杂
\makeatletter
\ifwhut@doctor
    \cdegree{工学博士}
\else
    \cdegree{工学硕士}
\fi
\makeatother

\cauthordepartment[计算机]{计算机科学与技术学院}
\cmajor{计算机科学与技术}
\cauthor{胡卫谊} 
\cadvisorname{陈天煌}
\cadvisortitle{副教授}
\cadvisordepartment{计算机科学与技术学院}
\cadvisordegree{本科}
\advisorzipcode{430063}
% 如果没有副指导老师，把下面几行相应的删除即可
\ccoadvisorname{林泓}
\ccoadvisortitle{副教授}
\ccoadvisordepartment{计算机科学与技术学院}
\ccoadvisordegree{本科}
\coadvisorzipcode{430063}
\cmajor{计算机科学与技术}
\csubmitdate{\the\year 年\the\month 月}
\coraldefensedate{\the\year 年\the\month 月}
\cauthorizedepartment{武汉理工大学}
\cauthorizedate{\the\year 年\the\month 月}
\cchairman{胡卫谊}
\cannotatora{胡卫谊a}
\cannotatorb{胡卫谊b}
%\cdate{\CTEXdigits{\the\year}年\CTEXnumber{\the\month}月}

\etitle{An Introduction to \LaTeX{} Thesis Template of Wuhan University of Technology}
\etitleformated{An Introduction to \LaTeX{} Thesis Template of\hfill \\ \hspace*{\fill} \hspace{9em}Wuhan University of Technology} 
% \edegree{Doctor of Science} 
\edegree{Doctor of Engineering} 
\emajor{Computer Science and Technology} 
\eauthor{Hu Weiyi} 
\eadvisorname{Professor Chen Tianhuang} 
\ecoadvisorname{Professor Lin Hong}
% 这个日期也会自动生成，你要改么？
% \edate{December, 2005}
\makecover
%% 中英文摘要
% 定义中英文摘要和关键字
\begin{cabstract}
  论文的摘要是对论文研究内容和成果的高度概括。摘要应对论文所研究的问题及其研究目
  的进行描述，对研究方法和过程进行简单介绍，对研究成果和所得结论进行概括。摘要应
  具有独立性和自明性，其内容应包含与论文全文同等量的主要信息。使读者即使不阅读全
  文，通过摘要就能了解论文的总体内容和主要成果。

  论文摘要的书写应力求精确、简明。切忌写成对论文书写内容进行提要的形式，尤其要避
  免“第 1 章……；第 2 章……；……”这种或类似的陈述方式。

  本文介绍武汉理工大学论文模板 \whutthesis{} 的使用方法。本模板符合学校的硕士和
  博士论文格式要求。

  本文的创新点主要有：
  \begin{itemize}
    \item 用例子来解释模板的使用方法；
    \item 用废话来填充无关紧要的部分；
    \item 一边学习摸索一边编写新代码。
  \end{itemize}

  关键词是为了文献标引工作、用以表示全文主要内容信息的单词或术语。关键词不超过 5
  个，每个关键词中间用分号分隔。（模板作者注：关键词分隔符不用考虑，模板会自动处
  理。英文关键词同理。）
\end{cabstract}

\ckeywords{\TeX, \LaTeX, CJK, 模板, 论文}

\begin{eabstract} 
   An abstract of a dissertation is a summary and extraction of research work
   and contributions. Included in an abstract should be description of research
   topic and research objective, brief introduction to methodology and research
   process, and summarization of conclusion and contributions of the
   research. An abstract should be characterized by independence and clarity and
   carry identical information with the dissertation. It should be such that the
   general idea and major contributions of the dissertation are conveyed without
   reading the dissertation. 

   An abstract should be concise and to the point. It is a misunderstanding to
   make an abstract an outline of the dissertation and words ``the first
   chapter'', ``the second chapter'' and the like should be avoided in the
   abstract.

   Key words are terms used in a dissertation for indexing, reflecting core
   information of the dissertation. An abstract may contain a maximum of 5 key
   words, with semi-colons used in between to separate one another.
\end{eabstract}

\ekeywords{\TeX, \LaTeX, CJK, template, thesis}

%% 目录
%%\cleardoublepage    %%这样页码显示是对的，但是pdf的书签指示的页码还是不对
%%\addcontentsline{toc}{chapter}{目录} % 将目录自身加入目录列表
\tableofcontents
%% 符号表
% !TeX root = ../main.tex

\begin{denotation}[3cm]
\item[HPC] 高性能计算 (High Performance Computing)
\item[cluster] 集群
\item[Itanium] 安腾
\item[SMP] 对称多处理
\item[API] 应用程序编程接口
\item[PI] 聚酰亚胺
\item[MPI] 聚酰亚胺模型化合物，N-苯基邻苯酰亚胺
\item[PBI] 聚苯并咪唑
\item[MPBI] 聚苯并咪唑模型化合物，N-苯基苯并咪唑
\item[PY] 聚吡咙
\item[PMDA-BDA]	均苯四酸二酐与联苯四胺合成的聚吡咙薄膜
\item[$\increment G$] 活化自由能 (Activation Free Energy)
\item[$\chi$] 传输系数 (Transmission Coefficient)
\item[$E$] 能量
\item[$m$] 质量
\item[$c$] 光速
\item[$P$] 概率
\item[$T$] 时间
\item[$v$] 速度
\item[劝学] 君子曰：学不可以已。青，取之于蓝，而青于蓝；冰，水为之，而寒于水。木
  直中绳。輮以为轮，其曲中规。虽有槁暴，不复挺者，輮使之然也。故木受绳则直，金就
  砺则利，君子博学而日参省乎己，则知明而行无过矣。吾尝终日而思矣，不如须臾之所学
  也；吾尝跂而望矣，不如登高之博见也。登高而招，臂非加长也，而见者远；顺风而呼，
  声非加疾也，而闻者彰。假舆马者，非利足也，而致千里；假舟楫者，非能水也，而绝江
  河，君子生非异也，善假于物也。积土成山，风雨兴焉；积水成渊，蛟龙生焉；积善成德，
  而神明自得，圣心备焉。故不积跬步，无以至千里；不积小流，无以成江海。骐骥一跃，
  不能十步；驽马十驾，功在不舍。锲而舍之，朽木不折；锲而不舍，金石可镂。蚓无爪牙
  之利，筋骨之强，上食埃土，下饮黄泉，用心一也。蟹六跪而二螯，非蛇鳝之穴无可寄托
  者，用心躁也。—— 荀况
\end{denotation}



% % 也可以使用 nomencl 宏包：

% \printnomenclature[3cm]

% \nomenclature{HPC}{高性能计算 (High Performance Computing)}
% \nomenclature{cluster}{集群}
% \nomenclature{Itanium}{安腾}
% \nomenclature{SMP}{对称多处理}
% \nomenclature{API}{应用程序编程接口}
% \nomenclature{PI}{聚酰亚胺}
% \nomenclature{MPI}{聚酰亚胺模型化合物，N-苯基邻苯酰亚胺}
% \nomenclature{PBI}{聚苯并咪唑}
% \nomenclature{MPBI}{聚苯并咪唑模型化合物，N-苯基苯并咪唑}
% \nomenclature{PY}{聚吡咙}
% \nomenclature{PMDA-BDA}{均苯四酸二酐与联苯四胺合成的聚吡咙薄膜}
% \nomenclature{$\increment G$}{活化自由能 (Activation Free Energy)}
% \nomenclature{$\chi$}{传输系数 (Transmission Coefficient)}
% \nomenclature{$E$}{能量}
% \nomenclature{$m$}{质量}
% \nomenclature{$c$}{光速}
% \nomenclature{$P$}{概率}
% \nomenclature{$T$}{时间}
% \nomenclature{$v$}{速度}

%%
%%
%% 正文
\mainmatter

%%% Local Variables:
%%% mode: latex
%%% TeX-master: t
%%% End:

\chapter{带 English 的标题}
\label{cha:intro}

这是 \whutthesis{} 的示例文档，基本上覆盖了模板中所有格式的设置。建议大家在使用模
板之前，除了阅读《\whutthesis{}用户手册》，这个示例文档也最好能看一看。

\mynote{这个对子有点意思。天对地，雨对风，大陆对长空，山花对海树，赤日对苍穹，老鼠的儿子会打洞……}
{小老鼠偷吃热凉粉；短长虫环绕矮高粱。}\footnote{韩愈（768-824），字退之，河南河阳（
  今河南孟县）人，自称郡望昌黎，世称韩昌黎。幼孤贫刻苦好学，德宗贞元八年进士。曾
  任监察御史，因上疏请免关中赋役，贬为阳山县令。后随宰相裴度平定淮西迁刑部侍郎，
  又因上表谏迎佛骨，贬潮州刺史。做过吏部侍郎，死谥文公，故世称韩吏部、韩文公。是
  唐代古文运动领袖，与柳宗元合称韩柳。诗力求险怪新奇，雄浑重气势。}


\section{封面相关}
封面的例子请参看 cover.tex。主要符号表参看 denation.tex，附录和个人简历分别参看 appendix01.tex
和 resume.tex。里面的命令都非常简单，一看即会。\footnote{你说还是看不懂？怎么会呢？}

\section{字体命令}
\label{sec:first}

苏轼（1037-1101），北宋文学家、书画家。字子瞻，号东坡居士，眉州眉山（今属四川）人
。苏洵子。嘉佑进士。神宗时曾任祠部员外郎，因反对王安石新法而求外职，任杭州通判，
知密州、徐州、湖州。后以作诗“谤讪朝廷”罪贬黄州。哲宗时任翰林学士，曾出知杭州、
颖州等，官至礼部尚书。后又贬谪惠州、儋州。北还后第二年病死常州。南宋时追谥文忠。
与父洵弟辙，合称“三苏”。在政治上属于旧党，但也有改革弊政的要求。其文汪洋恣肆，
明白畅达，为\mynote{中学语文经常考的，“韩柳与欧阳，三苏并曾王”，还记得不？}{“唐宋八大家”}之一。  其诗清新豪健，善用夸张比喻，在艺术表现方面独具
风格。少数诗篇也能反映民间疾苦，指责统治者的奢侈骄纵。词开豪放一派，对后代很有影
响。《念奴娇·赤壁怀古》、《水调歌头·丙辰中秋》传诵甚广。

{\kaishu 坡仙擅长行书、楷书，取法李邕、徐浩、颜真卿、杨凝式，而能自创新意。用笔丰腴
  跌宕，有天真烂漫之趣。与蔡襄、黄庭坚、米芾并称“宋四家”。能画竹，学文同，也喜
  作枯木怪石。论画主张“神似”，认为“论画以形似，见与儿童邻”；高度评价“诗中有
  画，画中有诗”的艺术造诣。诗文有《东坡七集》等。存世书迹有《答谢民师论文帖》、
  《祭黄几道文》、《前赤壁赋》、《黄州寒食诗帖》等。  画迹有《枯木怪石图》、《
  竹石图》等。}

{\fangsong 易与天地准，故能弥纶天地之道。仰以观於天文，俯以察於地理，是故知幽明之故。原
  始反终，故知死生之说。精气为物，游魂为变，是故知鬼神之情状。与天地相似，故不违。
  知周乎万物，而道济天下，故不过。旁行而不流，乐天知命，故不忧。安土敦乎仁，故
  能爱。范围天地之化而不过，曲成万物而不遗，通乎昼夜之道而知，故神无方而易无体。}

{\youyuan 有天地，然后万物生焉。盈天地之间者，唯万物，故受之以屯；屯者盈也，屯者物之
  始生也。物生必蒙，故受之以蒙；蒙者蒙也，物之穉也。物穉不可不养也，故受之以需；
  需者饮食之道也。饮食必有讼，故受之以讼。讼必有众起，故受之以师；师者众也。众必
  有所比，故受之以比；比者比也。比必有所畜也，故受之以小畜。物畜然后有礼，故受之
  以履。}

{\heiti 履而泰，然后安，故受之以泰；泰者通也。物不可以终通，故受之以否。物不可以终
  否，故受之以同人。与人同者，物必归焉，故受之以大有。有大者不可以盈，故受之以谦。
  有大而能谦，必豫，故受之以豫。豫必有随，故受之以随。以喜随人者，必有事，故受
  之以蛊；蛊者事也。}

{\lishu 有事而后可大，故受之以临；临者大也。物大然后可观，故受之以观。可观而后有所合
  ，故受之以噬嗑；嗑者合也。物不可以苟合而已，故受之以贲；贲者饰也。致饰然后亨
  ，则尽矣，故受之以剥；剥者剥也。物不可以终尽，剥穷上反下，故受之以复。复则不
  妄矣，故受之以无妄。}

{\songti 有无妄然后可畜，故受之以大畜。物畜然后可养，故受之以颐；颐者养也。不养则不
  可动，故受之以大过。物不可以终过，故受之以坎；坎者陷也。陷必有所丽，故受之以
  离；离者丽也。}

\section{测试行距的全文本页}
有无妄然后可畜，故受之以大畜。物畜然后可养，故受之以颐；颐者养也。不养则不
  可动，故受之以大过。物不可以终过，故受之以坎；坎者陷也。陷必有所丽，故受之以
  离；离者丽也。\\
  有无妄然后可畜，故受之以大畜。物畜然后可养，故受之以颐；颐者养也。不养则不
  可动，故受之以大过。物不可以终过，故受之以坎；坎者陷也。陷必有所丽，故受之以
  离；离者丽也。\\
  有无妄然后可畜，故受之以大畜。物畜然后可养，故受之以颐；颐者养也。不养则不
  可动，故受之以大过。物不可以终过，故受之以坎；坎者陷也。陷必有所丽，故受之以
  离；离者丽也。\\
  有无妄然后可畜，故受之以大畜。物畜然后可养，故受之以颐；颐者养也。不养则不
  可动，故受之以大过。物不可以终过，故受之以坎；坎者陷也。陷必有所丽，故受之以
  离；离者丽也。\\
  有无妄然后可畜，故受之以大畜。物畜然后可养，故受之以颐；颐者养也。不养则不
  可动，故受之以大过。物不可以终过，故受之以坎；坎者陷也。陷必有所丽，故受之以
  离；离者丽也。\\
  有无妄然后可畜，故受之以大畜。物畜然后可养，故受之以颐；颐者养也。不养则不
  可动，故受之以大过。物不可以终过，故受之以坎；坎者陷也。陷必有所丽，故受之以
  离；离者丽也。\\
  有无妄然后可畜，故受之以大畜。物畜然后可养，故受之以颐；颐者养也。不养则不
  可动，故受之以大过。物不可以终过，故受之以坎；坎者陷也。陷必有所丽，故受之以
  离；离者丽也。\\
  有无妄然后可畜，故受之以大畜。物畜然后可养，故受之以颐；颐者养也。不养则不
  可动，故受之以大过。物不可以终过，故受之以坎；坎者陷也。陷必有所丽，故受之以
  离；离者丽也。\\
  有无妄然后可畜，故受之以大畜。物畜然后可养，故受之以颐；颐者养也。不养则不
  可动，故受之以大过。物不可以终过，故受之以坎；坎者陷也。陷必有所丽，故受之以
  离；离者丽也。\\
  有无妄然后可畜，故受之以大畜。物畜然后可养，故受之以颐；颐者养也。不养则不
  可动，故受之以大过。物不可以终过，故受之以坎；坎者陷也。陷必有所丽，故受之以
  离；离者丽也。\\
  有无妄然后可畜，故受之以大畜。物畜然后可养，故受之以颐；颐者养也。不养则不
  可动，故受之以大过。物不可以终过，故受之以坎；坎者陷也。陷必有所丽，故受之以
  离；离者丽也。\\
  有无妄然后可畜，故受之以大畜。物畜然后可养，故受之以颐；颐者养也。不养则不
  可动，故受之以大过。物不可以终过，故受之以坎；坎者陷也。陷必有所丽，故受之以
  离；离者丽也。\\
  有无妄然后可畜，故受之以大畜。物畜然后可养，故受之以颐；颐者养也。不养则不
  可动，故受之以大过。物不可以终过，故受之以坎；坎者陷也。陷必有所丽，故受之以
  离；离者丽也。\\
  有无妄然后可畜，故受之以大畜。物畜然后可养，故受之以颐；颐者养也。不养则不
  可动，故受之以大过。物不可以终过，故受之以坎；坎者陷也。陷必有所丽，故受之以
  离；离者丽也。\\
  有无妄然后可畜，故受之以大畜。物畜然后可养，故受之以颐；颐者养也。不养则不
  可动，故受之以大过。物不可以终过，故受之以坎；坎者陷也。陷必有所丽，故受之以
  离；离者丽也。\\
  有无妄然后可畜，故受之以大畜。物畜然后可养，故受之以颐；颐者养也。不养则不
  可动，故受之以大过。物不可以终过，故受之以坎；坎者陷也。陷必有所丽，故受之以
  离；离者丽也。\\
  有无妄然后可畜，故受之以大畜。物畜然后可养，故受之以颐；颐者养也。不养则不
  可动，故受之以大过。物不可以终过，故受之以坎；坎者陷也。陷必有所丽，故受之以
  离；离者丽也。\\
\section{表格样本}
\label{chap1:sample:table} 

\subsection{基本表格}
\label{sec:basictable}

模板中关于表格的宏包有三个： \textsf{booktabs}、\textsf{array} 和
\textsf{longtabular}，命令有一个 \verb|\hlinewd|。三线表可以用 \textsf{booktabs}
提供的 \verb|\toprule|、\verb|\midrule| 和 \verb|\bottomrule|。它们与
\textsf{longtable} 能很好的配合使用。如果表格比较简单的话可以直接用命令
\verb|hlinewd{xpt}| 控制。
\begin{table}[htb]
  \centering
  \begin{minipage}[t]{0.8\linewidth} % 如果想在表格中使用脚注，minipage是个不错的办法
  \caption[模板文件]{模板文件。如果表格的标题很长，那么在表格索引中就会很不美
    观，所以要像 chapter 那样在前面用中括号写一个简短的标题。这个标题会出现在索
    引中。}
  \label{tab:template-files}
    \begin{tabular*}{\linewidth}{lp{10cm}}
      \toprule[1.5pt]
      {\heiti 文件名} & {\heiti 描述} \\\midrule[1pt]
      thuthesis.ins & \LaTeX{} 安装文件，docstrip\footnote{表格中的脚注} \\
      thuthesis.dtx & 所有的一切都在这里面\footnote{再来一个}。\\
      thuthesis.cls & 模板类文件。\\
      thuthesis.cfg & 模板配置文。cls 和 cfg 由前两个文件生成。\\
      thubib.bst    & 参考文献 Bibtex 样式文件。\\
      thutils.sty   & 常用的包和命令写在这里，减轻主文件的负担。\\
      \bottomrule[1.5pt]
    \end{tabular*}
  \end{minipage}
\end{table}

首先来看一个最简单的表格。表 \ref{tab:template-files} 列举了本模板主要文件及其功
能。请大家注意三线表中各条线对应的命令。这个例子还展示了如何在表格中正确使用脚注。
由于 \LaTeX{} 本身不支持在表格中使用 \verb|\footnote|，所以我们不得不将表格放在
小页中，而且最好将表格的宽度设置为小页的宽度，这样脚注看起来才更美观。

\subsection{复杂表格}
\label{sec:complicatedtable}

我们经常会在表格下方标注数据来源，或者对表格里面的条目进行解释。前面的脚注是一种
不错的方法，如果你不喜欢脚注。那么完全可以在表格后面自己写注释，比如表~\ref{tab:tabexamp1}。
\begin{table}[h]
  \centering
  \caption{复杂表格示例 1}  
  \label{tab:tabexamp1}
  \begin{minipage}[t]{0.8\textwidth} 
    \begin{tabularx}{\linewidth}{|l|X|X|X|X|}
      \hline
 \multirow{2}*{\backslashbox{x}{y}}  & \multicolumn{2}{c|}{First Half} & \multicolumn{2}{c|}{Second Half}\\\cline{2-5}
      & 1st Qtr &2nd Qtr&3rd Qtr&4th Qtr \\ \hline
      East$^{*}$ &   20.4&   27.4&   90&     20.4 \\
      West$^{**}$ &   30.6 &   38.6 &   34.6 &  31.6 \\ \hline
    \end{tabularx}\\[2pt]
    \footnotesize 注：数据来源《\whutthesis{} 使用手册》。\\
    *：东部\\
    **：西部
  \end{minipage}
\end{table}

此外，表~\ref{tab:tabexamp1} 同时还演示了另外两个功能：1）通过 \textsf{tabularx} 的
 \texttt{|X|} 扩展实现表格自动放大；2）通过命令 \verb|\backslashbox| 在表头部分
插入反斜线。

为了使我们的例子更接近实际情况，我会在必要的时候插入一些“无关”文字，以免太多图
表同时出现，导致排版效果不太理想。第一个出场的当然是我的最爱：风流潇洒、骏马绝尘、
健笔凌云的{\heiti 李太白}了。

李白，字太白，陇西成纪人。凉武昭王暠九世孙。或曰山东人，或曰蜀人。白少有逸才，志
气宏放，飘然有超世之心。初隐岷山，益州长史苏颋见而异之，曰：“是子天才英特，可比
相如。”天宝初，至长安，往见贺知章。知章见其文，叹曰：“子谪仙人也。”言于明皇，
召见金銮殿，奏颂一篇。帝赐食，亲为调羹，有诏供奉翰林。白犹与酒徒饮于市，帝坐沉香
亭子，意有所感，欲得白为乐章，召入，而白已醉。左右以水颒面，稍解，援笔成文，婉丽
精切。帝爱其才，数宴见。白常侍帝，醉，使高力士脱靴。力士素贵，耻之，摘其诗以激杨
贵妃。帝欲官白，妃辄沮止。白自知不为亲近所容，恳求还山。帝赐金放还。乃浪迹江湖，
终日沉饮。永王璘都督江陵，辟为僚佐。璘谋乱，兵败，白坐长流夜郎，会赦得还。族人阳
冰为当涂令，白往依之。代宗立，以左拾遗召，而白已卒。文宗时，诏以白歌诗、裴旻剑舞、
张旭草书为三绝云。集三十卷。今编诗二十五卷。\hfill\pozhehao《全唐诗》诗人小传

浮动体的并排放置一般有两种情况：1）二者没有关系，为两个独立的浮动体；2）二者隶属
于同一个浮动体。对表格来说并排表格既可以像图~\ref{tab:parallel1}、图~\ref{tab:parallel2} 
使用小页环境，也可以如图~\ref{tab:subtable} 使用子表格来做。图的例子参见第~\ref{sec:multifig} 节。
\begin{table}[h]
\noindent\begin{minipage}{0.5\textwidth}
\centering
\caption{第一个并排子表格}
\label{tab:parallel1}
\begin{tabular}{p{2cm}p{2cm}}
\toprule[1.5pt]
111 & 222 \\\midrule[1pt]
222 & 333 \\\bottomrule[1.5pt]
\end{tabular}
\end{minipage}
\begin{minipage}{0.5\textwidth}
\centering
\caption{第二个并排子表格}
\label{tab:parallel2}
\begin{tabular}{p{2cm}p{2cm}}
\toprule[1.5pt]
111 & 222 \\\midrule[1pt]
222 & 333 \\\bottomrule[1.5pt]
\end{tabular}
\end{minipage}
\end{table}

然后就是忧国忧民，诗家楷模杜工部了。杜甫，字子美，其先襄阳人，曾祖依艺为巩令，因
居巩。甫天宝初应进士，不第。后献《三大礼赋》，明皇奇之，召试文章，授京兆府兵曹参
军。安禄山陷京师，肃宗即位灵武，甫自贼中遁赴行在，拜左拾遗。以论救房琯，出为华州
司功参军。关辅饥乱，寓居同州同谷县，身自负薪采梠，餔糒不给。久之，召补京兆府功曹，
道阻不赴。严武镇成都，奏为参谋、检校工部员外郎，赐绯。武与甫世旧，待遇甚厚。乃于
成都浣花里种竹植树，枕江结庐，纵酒啸歌其中。武卒，甫无所依，乃之东蜀就高適。既至
而適卒。是岁，蜀帅相攻杀，蜀大扰。甫携家避乱荆楚，扁舟下峡，未维舟而江陵亦乱。乃
溯沿湘流，游衡山，寓居耒阳。卒年五十九。元和中，归葬偃师首阳山，元稹志其墓。天宝
间，甫与李白齐名，时称李杜。然元稹之言曰：“李白壮浪纵恣，摆去拘束，诚亦差肩子美
矣。至若铺陈终始，排比声韵，大或千言，次犹数百，词气豪迈，而风调清深，属对律切，
而脱弃凡近，则李尚不能历其藩翰，况堂奥乎。”白居易亦云：“杜诗贯穿古今，  尽工尽
善，殆过于李。”元、白之论如此。盖其出处劳佚，喜乐悲愤，好贤恶恶，一见之于诗。而
又以忠君忧国、伤时念乱为本旨。读其诗可以知其世，故当时谓之“诗史”。旧集诗文共六
十卷，今编诗十九卷。

\begin{table}
\centering
\caption{并排子表格}
\label{tab:subtable}
\subfloat[第一个子表格]{
\begin{tabular}{p{2cm}p{2cm}}
\toprule[1.5pt]
111 & 222 \\\midrule[1pt]
222 & 333 \\\bottomrule[1.5pt]
\end{tabular}}\hskip2cm
\subfloat[第二个子表格]{
\begin{tabular}{p{2cm}p{2cm}}
\toprule[1.5pt]
111 & 222 \\\midrule[1pt]
222 & 333 \\\bottomrule[1.5pt]
\end{tabular}}
\end{table}

不可否认 \LaTeX{} 的表格功能没有想象中的那么强大，不过只要你足够认真，足够细致，那么
同样可以排出来非常复杂非常漂亮的表格。请参看表~\ref{tab:tabexamp2}。
\begin{table}[hb]
  \centering\zihao{5}
  \caption{复杂表格示例 2}
  \label{tab:tabexamp2}
  \begin{tabular}[c]{|c|m{0.8in}|c|c|c|c|c|}\hline
    \multicolumn{2}{|c|}{Network Topology} & \# of nodes & 
    \multicolumn{3}{c|}{\# of clients} & Server \\\hline
    GT-ITM & Waxman Transit-Stub & 600 &
    \multirow{2}{2em}{2\%}& 
    \multirow{2}{2em}{10\%}& 
    \multirow{2}{2em}{50\%}& 
    \multirow{2}{1.2in}{Max. Connectivity}\\\cline{1-3}
    \multicolumn{2}{|c|}{Inet-2.1} & 6000 & & & &\\\hline
    \multirow{2}{1in}{Xue} & Rui  & Ni &\multicolumn{4}{c|}{\multirow{2}*{\whutthesis}}\\\cline{2-3}
    & \multicolumn{2}{c|}{ABCDEF} &\multicolumn{4}{c|}{} \\\hline
\end{tabular}
\end{table}

最后就是清新飘逸、文约意赅、空谷绝响的王大侠了。王维，字摩诘，河东人。工书画，与
弟缙俱有俊才。开元九年，进士擢第，调太乐丞。坐累为济州司仓参军，历右拾遗、监察御
史、左补阙、库部郎中，拜吏部郎中。天宝末，为给事中。安禄山陷两都，维为贼所得，服
药阳喑，拘于菩提寺。禄山宴凝碧池，维潜赋诗悲悼，闻于行在。贼平，陷贼官三等定罪，
特原之，责授太子中允，迁中庶子、中书舍人。复拜给事中，转尚书右丞。维以诗名盛于开
元、天宝间，宁薛诸王驸马豪贵之门，无不拂席迎之。得宋之问辋川别墅，山水绝胜，与道
友裴迪，浮舟往来，弹琴赋诗，啸咏终日。笃于奉佛，晚年长斋禅诵。一日，忽索笔作书
数纸，别弟缙及平生亲故，舍笔而卒。赠秘书监。宝应中，代宗问缙：“朕常于诸王坐闻维
乐章，今存几何？”缙集诗六卷，文四卷，表上之。敕答云，卿伯氏位列先朝，名高希代。
抗行周雅，长揖楚辞。诗家者流，时论归美。克成编录，叹息良深。殷璠谓维诗词秀调雅，
意新理惬。在泉成珠，著壁成绘。苏轼亦云：“维诗中有画，画中有诗也。”今编诗四卷。

要想用好论文模板还是得提前学习一些 \TeX/\LaTeX{}的相关知识，具备一些基本能力，掌
握一些常见技巧，否则一旦遇到问题还真是比较麻烦。我们见过很多这样的同学，一直以来
都是使用 Word 等字处理工具，以为 \LaTeX{}模板的用法也应该类似，所以就沿袭同样的思
路来对待这种所见非所得的排版工具，结果被折腾的焦头烂额，疲惫不堪。

如果您要排版的表格长度超过一页，那么推荐使用 \textsf{longtable} 或者 \textsf{supertabular} 
宏包，模板对 \textsf{longtable} 进行了相应的设置，所以用起来可能简单一些。
表~\ref{tab:performance} 就是 \textsf{longtable} 的简单示例。
\begin{longtable}[c]{c*{6}{r}}
\caption{实验数据}\label{tab:performance}\\
\toprule[1.5pt]
 测试程序 & \multicolumn{1}{c}{正常运行} & \multicolumn{1}{c}{同步} & \multicolumn{1}{c}{检查点} & \multicolumn{1}{c}{卷回恢复}
& \multicolumn{1}{c}{进程迁移} & \multicolumn{1}{c}{检查点} \\
& \multicolumn{1}{c}{时间 (s)}& \multicolumn{1}{c}{时间 (s)}&
\multicolumn{1}{c}{时间 (s)}& \multicolumn{1}{c}{时间 (s)}& \multicolumn{1}{c}{
  时间 (s)}&  文件（KB）\\\midrule[1pt]
\endfirsthead
\multicolumn{7}{c}{续表~\thetable\hskip1em 实验数据}\\
\toprule[1.5pt]
 测试程序 & \multicolumn{1}{c}{正常运行} & \multicolumn{1}{c}{同步} & \multicolumn{1}{c}{检查点} & \multicolumn{1}{c}{卷回恢复}
& \multicolumn{1}{c}{进程迁移} & \multicolumn{1}{c}{检查点} \\
& \multicolumn{1}{c}{时间 (s)}& \multicolumn{1}{c}{时间 (s)}&
\multicolumn{1}{c}{时间 (s)}& \multicolumn{1}{c}{时间 (s)}& \multicolumn{1}{c}{
  时间 (s)}&  文件（KB）\\\midrule[1pt]
\endhead
\hline
\multicolumn{7}{r}{续下页}
\endfoot
\endlastfoot
CG.A.2 & 23.05 & 0.002 & 0.116 & 0.035 & 0.589 & 32491 \\
CG.A.4 & 15.06 & 0.003 & 0.067 & 0.021 & 0.351 & 18211 \\
CG.A.8 & 13.38 & 0.004 & 0.072 & 0.023 & 0.210 & 9890 \\
CG.B.2 & 867.45 & 0.002 & 0.864 & 0.232 & 3.256 & 228562 \\
CG.B.4 & 501.61 & 0.003 & 0.438 & 0.136 & 2.075 & 123862 \\
CG.B.8 & 384.65 & 0.004 & 0.457 & 0.108 & 1.235 & 63777 \\
MG.A.2 & 112.27 & 0.002 & 0.846 & 0.237 & 3.930 & 236473 \\
MG.A.4 & 59.84 & 0.003 & 0.442 & 0.128 & 2.070 & 123875 \\
MG.A.8 & 31.38 & 0.003 & 0.476 & 0.114 & 1.041 & 60627 \\
MG.B.2 & 526.28 & 0.002 & 0.821 & 0.238 & 4.176 & 236635 \\
MG.B.4 & 280.11 & 0.003 & 0.432 & 0.130 & 1.706 & 123793 \\
MG.B.8 & 148.29 & 0.003 & 0.442 & 0.116 & 0.893 & 60600 \\
LU.A.2 & 2116.54 & 0.002 & 0.110 & 0.030 & 0.532 & 28754 \\
LU.A.4 & 1102.50 & 0.002 & 0.069 & 0.017 & 0.255 & 14915 \\
LU.A.8 & 574.47 & 0.003 & 0.067 & 0.016 & 0.192 & 8655 \\
LU.B.2 & 9712.87 & 0.002 & 0.357 & 0.104 & 1.734 & 101975 \\
LU.B.4 & 4757.80 & 0.003 & 0.190 & 0.056 & 0.808 & 53522 \\
LU.B.8 & 2444.05 & 0.004 & 0.222 & 0.057 & 0.548 & 30134 \\
EP.A.2 & 123.81 & 0.002 & 0.010 & 0.003 & 0.074 & 1834 \\
EP.A.4 & 61.92 & 0.003 & 0.011 & 0.004 & 0.073 & 1743 \\
EP.A.8 & 31.06 & 0.004 & 0.017 & 0.005 & 0.073 & 1661 \\
EP.B.2 & 495.49 & 0.001 & 0.009 & 0.003 & 0.196 & 2011 \\
EP.B.4 & 247.69 & 0.002 & 0.012 & 0.004 & 0.122 & 1663 \\
EP.B.8 & 126.74 & 0.003 & 0.017 & 0.005 & 0.083 & 1656 \\
\bottomrule[1.5pt]
\end{longtable}

\subsection{其它}
\label{sec:tableother}
有的同学不想让某个表格或者图片出现在索引里面，那么请使用命令 \verb|\caption*{}|，
这个命令不会给表格编号，也就是出来的只有标题文字而没有“表~XX”，“图~XX”，否则
索引里面序号不连续就显得不伦不类，这也是 \LaTeX{} 里星号命令默认的规则。

有这种需求的多是本科同学的英文资料翻译部分，如果你觉得附录中英文原文中的表格和图
片显示成“  表”和“图”很不协调的话，一个很好的办法就是用 \verb|\caption*|，参数
随便自己写，比如不守规矩的表~1.111 和图~1.111 能满足这种特殊需要（可以参看附录部
分）。
\begin{table}[ht]
\centering
  \begin{minipage}{0.45\linewidth}
  \centering
  \caption*{表~1.111\hskip1em 这是一个手动编号，不出现在索引中的表格。}
  \label{tab:badtabular}
  \begin{picture}(150,50)
    \framebox(150,50)[c]{\whutthesis}
  \end{picture}    
  \end{minipage}\hfill
  \begin{minipage}{0.45\linewidth}
  \centering
  \begin{picture}(150,50)
    \framebox(150,50)[c]{薛瑞尼}
  \end{picture}
  \caption*{Figure~1.111\hskip1em 这是一个手动编号，不出现在索引中的图。}
  \label{tab:badfigure}
  \end{minipage}
\end{table}

如果你的确想让它编号，但又不想让它出现在索引中的话，那就自己看看代码改一改吧，我
目前不打算给模板增加这种另类命令。

最后，虽然大家不一定会独立使用小页，但是关于小页中的脚注还是有必要提一下。请看下
面的例子。

\begin{minipage}[t]{\linewidth-2\parindent}
  柳宗元，字子厚（773-819），河东（今永济县）人\footnote{山西永济水饺。}，是唐代
  杰出的文学家，哲学家，同时也是一位政治改革家。与韩愈共同倡导唐代古文运动，并称
  韩柳\footnote{唐宋八大家之首二位。}。
\end{minipage}\\[-5pt]

唐朝安史之乱后，宦官专权，藩镇割据，土地兼并日渐严重，社会生产破坏严重，民不聊生。柳宗
元对这种社会现实极为不满，他积极参加了王叔文领导的“永济革新”，并成为这一
运动的中坚人物。他们革除弊政，打击权奸，触犯了宦官和官僚贵族利益，在他们的联合反
扑下，改革失败了，柳宗元被贬为永州司马。

\section{定理环境}
\label{sec:theorem}

给大家演示一下各种和证明有关的环境：

\begin{assumption}
待月西厢下，迎风户半开；隔墙花影动，疑是玉人来。
\begin{eqnarray}
  \label{eq:eqnxmp}
  c & = & a^2 - b^2\\
    & = & (a+b)(a-b)
\end{eqnarray}
\end{assumption}

千辛万苦，历尽艰难，得有今日。然相从数千里，未曾哀戚。今将渡江，方图百年欢笑，如
何反起悲伤？（引自《杜十娘怒沉百宝箱》）

\begin{definition}
子曰：「道千乘之国，敬事而信，节用而爱人，使民以时。」
\end{definition}

千古第一定义！问世间、情为何物，只教生死相许？天南地北双飞客，老翅几回寒暑。欢乐趣，离别苦，就中更有痴儿女。
君应有语，渺万里层云，千山暮雪，只影向谁去？

横汾路，寂寞当年箫鼓，荒烟依旧平楚。招魂楚些何嗟及，山鬼暗谛风雨。天也妒，未信与，莺儿燕子俱黄土。
千秋万古，为留待骚人，狂歌痛饮，来访雁丘处。

\begin{proposition}
 曾子曰：「吾日三省吾身 \pozhehao 为人谋而不忠乎？与朋友交而不信乎？传不习乎？」
\end{proposition}

多么凄美的命题啊！其日牛马嘶，新妇入青庐，奄奄黄昏后，寂寂人定初，我命绝今日，
魂去尸长留，揽裙脱丝履，举身赴清池，府吏闻此事，心知长别离，徘徊庭树下，自挂东南
枝。

\begin{remark}
天不言自高，水不言自流。
\begin{gather*}
\begin{split} 
\varphi(x,z)
&=z-\gamma_{10}x-\gamma_{mn}x^mz^n\\
&=z-Mr^{-1}x-Mr^{-(m+n)}x^mz^n
\end{split}\\[6pt]
\begin{align} \zeta^0&=(\xi^0)^2,\\
\zeta^1 &=\xi^0\xi^1,\\
\zeta^2 &=(\xi^1)^2,
\end{align}
\end{gather*}
\end{remark}

天尊地卑，乾坤定矣。卑高以陈，贵贱位矣。 动静有常，刚柔断矣。方以类聚，物以群分，
吉凶生矣。在天成象，在地成形，变化见矣。鼓之以雷霆，润之以风雨，日月运行，一寒一
暑，乾道成男，坤道成女。乾知大始，坤作成物。乾以易知，坤以简能。易则易知，简则易
从。易知则有亲，易从则有功。有亲则可久，有功则可大。可久则贤人之德，可大则贤人之
业。易简，而天下矣之理矣；天下之理得，而成位乎其中矣。

\begin{axiom}
两点间直线段距离最短。  
\begin{align}
x&\equiv y+1\pmod{m^2}\\
x&\equiv y+1\mod{m^2}\\
x&\equiv y+1\pod{m^2}
\end{align}
\end{axiom}

《彖曰》：大哉乾元，万物资始，乃统天。云行雨施，品物流形。大明始终，六位时成，时
乘六龙以御天。乾道变化，各正性命，保合大和，乃利贞。首出庶物，万国咸宁。

《象曰》：天行健，君子以自强不息。潜龙勿用，阳在下也。见龙再田，德施普也。终日乾
乾，反复道也。或跃在渊，进无咎也。飞龙在天，大人造也。亢龙有悔，盈不可久也。用九，
天德不可为首也。 　　

\begin{lemma}
《猫和老鼠》是我最爱看的动画片。
\begin{multline*}%\tag*{[a]} % 这个不出现在索引中
\int_a^b\biggl\{\int_a^b[f(x)^2g(y)^2+f(y)^2g(x)^2]
 -2f(x)g(x)f(y)g(y)\,dx\biggr\}\,dy \\
 =\int_a^b\biggl\{g(y)^2\int_a^bf^2+f(y)^2
  \int_a^b g^2-2f(y)g(y)\int_a^b fg\biggr\}\,dy
\end{multline*}
\end{lemma}

行行重行行，与君生别离。相去万余里，各在天一涯。道路阻且长，会面安可知。胡马依北
风，越鸟巢南枝。相去日已远，衣带日已缓。浮云蔽白日，游子不顾返。思君令人老，岁月
忽已晚。  弃捐勿复道，努力加餐饭。

\begin{theorem}\label{the:theorem1}
犯我强汉者，虽远必诛\hfill \pozhehao 陈汤（汉）
\end{theorem}
\begin{subequations}
\begin{align}
y & = 1 \\
y & = 0
\end{align}
\end{subequations}
道可道，非常道。名可名，非常名。无名天地之始；有名万物之母。故常无，欲以观其妙；
常有，欲以观其徼。此两者，同出而异名，同谓之玄。玄之又玄，众妙之门。上善若水。水
善利万物而不争，处众人之所恶，故几于道。曲则全，枉则直，洼则盈，敝则新，少则多，
多则惑。人法地，地法天，天法道，道法自然。知人者智，自知者明。胜人者有力，自胜
者强。知足者富。强行者有志。不失其所者久。死而不亡者寿。

\begin{proof}
燕赵古称多感慨悲歌之士。董生举进士，连不得志于有司，怀抱利器，郁郁适兹土，吾
知其必有合也。董生勉乎哉？

夫以子之不遇时，苟慕义强仁者，皆爱惜焉，矧燕、赵之士出乎其性者哉！然吾尝闻
风俗与化移易，吾恶知其今不异于古所云邪？聊以吾子之行卜之也。董生勉乎哉？

吾因子有所感矣。为我吊望诸君之墓，而观于其市，复有昔时屠狗者乎？为我谢
曰：“明天子在上，可以出而仕矣！” \hfill\pozhehao 韩愈《送董邵南序》
\end{proof}

\begin{corollary}
  四川话配音的《猫和老鼠》是世界上最好看最好听最有趣的动画片。
\begin{alignat}{3}
V_i & =v_i - q_i v_j, & \qquad X_i & = x_i - q_i x_j,
 & \qquad U_i & = u_i,
 \qquad \text{for $i\ne j$;}\label{eq:B}\\
V_j & = v_j, & \qquad X_j & = x_j,
  & \qquad U_j & u_j + \sum_{i\ne j} q_i u_i.
\end{alignat}
\end{corollary}

迢迢牵牛星，皎皎河汉女。
纤纤擢素手，札札弄机杼。
终日不成章，泣涕零如雨。
河汉清且浅，相去复几许。
盈盈一水间，脉脉不得语。

\begin{example}
  大家来看这个例子。
\begin{equation}
\label{ktc}
\left\{\begin{array}{l}
\nabla f({\mbox{\boldmath $x$}}^*)-\sum\limits_{j=1}^p\lambda_j\nabla g_j({\mbox{\boldmath $x$}}^*)=0\\[0.3cm]
\lambda_jg_j({\mbox{\boldmath $x$}}^*)=0,\quad j=1,2,\cdots,p\\[0.2cm]
\lambda_j\ge 0,\quad j=1,2,\cdots,p.
\end{array}\right.
\end{equation}
\end{example}

\begin{exercise}
  清列出 Andrew S. Tanenbaum 和 W. Richard Stevens 的所有著作。
\end{exercise}

\begin{conjecture} \textit{Poincare Conjecture} If in a closed three-dimensional
  space, any closed curves can shrink to a point continuously, this space can be
  deformed to a sphere.
\end{conjecture}

\begin{problem}
 回答还是不回答，是个问题。 
\end{problem}

如何引用定理~\ref{the:theorem1} 呢？加上 \verb|label| 使用 \verb|ref| 即可。妾发
初覆额，折花门前剧。郎骑竹马来，绕床弄青梅。同居长干里，两小无嫌猜。 十四为君妇，
羞颜未尝开。低头向暗壁，千唤不一回。十五始展眉，愿同尘与灰。常存抱柱信，岂上望夫
台。 十六君远行，瞿塘滟滪堆。五月不可触，猿声天上哀。门前迟行迹，一一生绿苔。苔深
不能扫，落叶秋风早。八月蝴蝶来，双飞西园草。感此伤妾心，坐愁红颜老。

\section{参考文献}
\label{sec:bib}
当然参考文献可以直接写 bibitem，虽然费点功夫，但是好控制，各种格式可以自己随意改
写。

本模板推荐使用 BIB\TeX，样式文件为 GBT7714-2005NLang.bst，基本符合学校的参考文献格式（如专利
等引用未加详细测试）。看看这个例子，关于书的\cite{tex, companion, ColdSources}，
还有这些\cite{Krasnogor2004e, clzs, zjsw}，关于杂志的\cite{ELIDRISSI94,
  MELLINGER96, SHELL02}，硕士论文\cite{zhubajie, metamori2004}，博士论文
\cite{shaheshang, FistSystem01}，标准文件\cite{IEEE-1363}，会议论文\cite{DPMG,kocher99}，技术报告\cite{NPB2}。中文参
考文献\cite{cnarticle}应增加 \texttt{lang=``zh''} 字段，以便进行相应处理。另
外，这个 bst 对中文文献\cite{cnproceed}的支持并不是十全十美，如果有不如意的地方，
请手动修改 bbl 文件。

有时候不想要上标，那么可以这样 \onlinecite{shaheshang}，这个非常重要。

\section{公式}
\label{sec:equation}
贝叶斯公式如式~(\ref{equ:chap1:bayes})，其中 $p(y|\mathbf{x})$ 为后验；
$p(\mathbf{x})$ 为先验；分母 $p(\mathbf{x})$ 为归一化因子。
\begin{equation}
\label{equ:chap1:bayes}
p(y|\mathbf{x}) = \frac{p(\mathbf{x},y)}{p(\mathbf{x})}=
\frac{p(\mathbf{x}|y)p(y)}{p(\mathbf{x})} 
\end{equation}

论文里面公式越多，\TeX{} 就越 happy。再看一个 \textsf{amsmath} 的例子：
\newcommand{\envert}[1]{\left\lvert#1\right\rvert} 
\begin{equation}\label{detK2}
\det\mathbf{K}(t=1,t_1,\dots,t_n)=\sum_{I\in\mathbf{n}}(-1)^{\envert{I}}
\prod_{i\in I}t_i\prod_{j\in I}(D_j+\lambda_jt_j)\det\mathbf{A}
^{(\lambda)}(\overline{I}|\overline{I})=0.
\end{equation} 

前面定理示例部分列举了很多公式环境，可以说把常见的情况都覆盖了，大家在写公式的时
候一定要好好看 \textsf{amsmath} 的文档，并参考模板中的用法：
\begin{multline*}%\tag{[b]} % 这个出现在索引中的
\int_a^b\biggl\{\int_a^b[f(x)^2g(y)^2+f(y)^2g(x)^2]
 -2f(x)g(x)f(y)g(y)\,dx\biggr\}\,dy \\
 =\int_a^b\biggl\{g(y)^2\int_a^bf^2+f(y)^2
  \int_a^b g^2-2f(y)g(y)\int_a^b fg\biggr\}\,dy
\end{multline*}

其实还可以看看这个多级规划：
\begin{equation}\label{bilevel}
\left\{\begin{array}{l}
\max\limits_{{\mbox{\footnotesize\boldmath $x$}}} F(x,y_1^*,y_2^*,\cdots,y_m^*)\\[0.2cm]
\mbox{subject to:}\\[0.1cm]
\qquad G(x)\le 0\\[0.1cm]
\qquad(y_1^*,y_2^*,\cdots,y_m^*)\mbox{ solves problems }(i=1,2,\cdots,m)\\[0.1cm]
\qquad\left\{\begin{array}{l}
    \max\limits_{{\mbox{\footnotesize\boldmath $y_i$}}}f_i(x,y_1,y_2,\cdots,y_m)\\[0.2cm]
    \mbox{subject to:}\\[0.1cm]
    \qquad g_i(x,y_1,y_2,\cdots,y_m)\le 0.
    \end{array}\right.
\end{array}\right.
\end{equation}
这些跟规划相关的公式都来自于刘宝碇老师《不确定规划》的课件。

\section{破折号}
\label{sec:pozhehao}

中文破折号为一个两个字宽垂直居中的直线，输入法直接得到的破折号是两个断开的小短线
（——），这看起来不舒服。所以我定义了一个破折号的命令 \verb|\pozhehao|，请看几个
例子：
\begin{itemize}
\item 这是一个 \pozhehao 破折号
  \begin{enumerate}[(1)]
  \item 同时也可以看看
  \item 不同列表环境的间距
  \end{enumerate}
\item 看起来这个要好一些
\item 破折号 \pozhehao 就说到这里。
\end{itemize}

默认的列表环境上下间距很大，模板将其重定义为 \textsf{paralist} 中的压缩环境，看起
来要好一些。如果还是不满意，自己也可以调 \verb|\itemsep| 的。\textsf{paralist} 还
可以方便的指定标签的样式。



%%% Local Variables: 
%%% mode: latex
%%% TeX-master: t
%%% End: 

\chapter{中华人民共和国}
\label{cha:china}

\section{其它例子}
\label{sec:other}

在第~\ref{cha:intro} 章中我们学习了贝叶斯公式~(\ref{equ:chap1:bayes})，这里我们复
习一下：
\begin{equation}
\label{equ:chap2:bayes}
p(y|\mathbf{x}) = \frac{p(\mathbf{x},y)}{p(\mathbf{x})}=
\frac{p(\mathbf{x}|y)p(y)}{p(\mathbf{x})} 
\end{equation}

\subsection{绘图}
\label{sec:draw}

本模板不再预先装载任何绘图包（如 \textsf{pstricks，pgf} 等），完全由你自己来决定。
个人觉得 \textsf{pgf} 不错，不依赖于 Postscript。此外还有很多针对 \LaTeX{} 的
 GUI 作图工具，如 XFig(jFig), WinFig, Tpx, Ipe, Dia, Inkscape, LaTeXPiX,
jPicEdt, jaxdraw 等等。

\subsection{插图}
\label{sec:graphs}

强烈推荐《\LaTeXe 插图指南》！关于子图形的使用细节请参看 \textsf{subfig} 的说明文档。

\subsubsection{一个图形}
\label{sec:onefig}
一般图形都是处在浮动环境中。之所以称为浮动是指最终排版效果图形的位置不一定与源文
件中的位置对应\footnote{This is not a bug, but a feature of \LaTeX!}，这也是刚使
用 \LaTeX{} 同学可能遇到的问题。如果要强制固定浮动图形的位置，请使用 \textsf{float} 宏包，
它提供了 \texttt{[H]} 参数，比如图~\ref{fig:xfig1}。
\begin{figure}[H] % use float package if you want it here
  \centering
  \includegraphics{hello}
  \caption{利用 Xfig 制图}
  \label{fig:xfig1}
\end{figure}

大学之道，在明明德，在亲民，在止于至善。知止而后有定；定而后能静；静而后能安；安
而后能虑；虑而后能得。物有本末，事有终始。知所先后，则近道矣。古之欲明明德于天
下者，先治其国；欲治其国者，先齐其家；欲齐其家者，先修其身；欲修其身者，先正其心；
欲正其心者，先诚其意；欲诚其意者，先致其知；致知在格物。物格而后知至；知至而后
意诚；意诚而后心正；心正而后身 修；身修而后家齐；家齐而后国治；国治而后天下
平。自天子以至于庶人，壹是皆以修身为本。其本乱而未治者 否矣。其所厚者薄，而其所
薄者厚，未之有也！

\hfill \pozhehao《大学》


\subsubsection{多个图形}
\label{sec:multifig}

如果多个图形相互独立，并不共用一个图形计数器，那么用 \verb|minipage| 或者
\verb|parbox| 就可以。否则，请参看图~\ref{fig:big1}，它包含两个小图，分别是图~\ref{fig:subfig1} 
和图~\ref{fig:subfig2}。推荐使用 \verb|\subfloat|，不要再用
\verb|\subfigure| 和 \verb|\subtable|。
\begin{figure} %[h]
  \centering%
  \subfloat[第一个小图形]{%
    \label{fig:subfig1}
    \includegraphics[height=2cm]{thu-fig-logo}}\hspace{4em}%
  \subfloat[第二个小图形。如果标题很长的话，它会自动换行，这个 caption 就是这样的例子]{%
    \label{fig:subfig2}
    \includegraphics[height=2cm]{thu-text-logo}}
  \caption{包含子图形的大图形}
  \label{fig:big1}
\end{figure}

古之学者必有师。师者，所以传道受业解惑也。人非生而知之者，孰能无惑？惑而不从师，
其为惑也，终不解矣。生乎吾前，其闻道也固先乎吾，吾从而师之；生乎吾後，其闻道也亦
先乎吾，吾从而师之。吾师道也，夫庸知其年之先後生於吾乎！是故无贵无贱无长无少，道
之所存，师之所存也。

嗟乎！师道之不传也久矣，欲人之无惑也难矣。古之圣人，其出人也远矣，犹且从师而问焉；
今之众人，其下圣人也亦远矣，而耻学於师。是故圣益圣，愚益愚。圣人之所以为圣，愚
人之所以为愚，其皆出於此乎？爱其子，择师而教之，於其身也，则耻师焉，惑焉。彼童子
之师，授之书而习其句读者，非吾所谓传其道、解其惑者也。句读之不知，惑之不解，或师
焉，或不焉，小学而大遗，吾未见其明也。巫医、乐师、百工之人不耻相师，  士大夫之族
曰“师”曰“弟子”之云者，则群聚而笑之。问之，则曰：彼与彼年相若也，道相似也，位
卑则足羞，官盛则近谀。呜呼！师道之不复，可知矣。巫医、乐师、百工之人。吾子不齿，
今其智乃反不能及，其可怪也欤！圣人无常师。孔子师郯子、苌子、师襄、老聃。郯子之徒，
其贤不及孔子。孔子曰：“三人行，必有我师。”是故弟子不必不如师，师不必贤於弟子。
闻道有先後，术业有专攻，如是而已。

如果要把编号的两个图形并排，那么小页就非常有用了：
\begin{figure}
\begin{minipage}{0.48\textwidth}
  \centering
  \includegraphics[height=2cm]{thu-whole-logo}
  \caption{并排第一个图}
  \label{fig:parallel1}
\end{minipage}\hfill
\begin{minipage}{0.48\textwidth}
  \centering
  \includegraphics[height=2cm]{thu-whole-logo}
  \caption{并排第二个图}
  \label{fig:parallel2}
\end{minipage}
\end{figure}

李氏子蟠，年十七，好古文、六艺，经传皆通习之，不拘於时，学於余。余嘉其能行古
道，作师说以贻之。

\hfill \pozhehao 韩愈（唐）


\section{用~listings~插入源代码}

原先~ctexbook~文档类和~listings~宏包配合使用时，代码在换页时会出现莫名其妙的错误，后来经高人指点，顺利解决了。
感兴趣的话，可以看看\href{http://bbs.ctex.org/viewthread.php?tid=53451}{这里}。
这里给使用~listings~宏包插入源代码的例子，这里是一段C代码。
另外，listings宏包真可谓博大精深，可以实现各种复杂、漂亮的效果，想要进一步学习的同学，可以参考
\href{http://mirror.ctan.org/macros/latex/contrib/listings/listings.pdf}{listings宏包手册}。

\begin{lstlisting}[language={C}, caption={一段C源代码}]
#include <stdio.h>
#include <unistd.h>
#include <sys/types.h>
#include <sys/wait.h>

int main() {
  pid_t pid;

  switch ((pid = fork())) {
  case -1:
    printf("fork failed\n");
    break;
  case 0:
    /* child calls exec */
    execl("/bin/ls", "ls", "-l", (char*)0);
    printf("execl failed\n");
    break;
  default:
    /* parent uses wait to suspend execution until child finishes */
    wait((int*)0);
    printf("is completed\n");
    break;
  }

  return 0;
}
\end{lstlisting}

再给一个插入MATLAB代码的例子，感谢~daisying~站友提供的代码。

\begin{lstlisting}[language={matlab}, caption={一段MATLAB源代码}]
function paper1
r=0.05;
n=100;
T=1;
X=1;
v0=0.8;
sigma=sqrt(0.08);
deltat=T/n;
for i=1:n
    t(i)=i*deltat;
    w(i)=random('norm',0,t(i),1);
end
for i=1:n
    alpha(i)=0.39;
end
for i=1:n
    temp=0;
    for k=1:i
        temp=temp+alpha(k);
    end
    B(i)=exp(r*t(i));
    BB(i)=B(i)*exp(temp*deltat);
    BBB(i)=exp(-r*(T-t(i)));
end
for i=1:n
    s0(i)=X*BBB(i);
    v(i)=v0*exp((r-0.5*sigma^2)*t(i)+sigma*w(i));
    for j=i+1:n
        D=X*BBB(j);
        d1=(log(v(i)/D)+(r+sigma^2/2)*(t(j)-t(i)))/(sigma*sqrt(t(j)-t(i)));
        d2=d1-(sigma*sqrt(t(j)-t(i)));
        ppp(i,j)=D*exp(-r*(t(j)-t(i)))*(1-cdf('normal',d2,0,1))-v(i)*(1-cdf('n
ormal',d1,0,1));
    end
end
for i=1:n
    s1(i)=0;
    for j=i+1:n
        s1(i)=s1(i)+BB(j)^(-1)*alpha(j)*deltat*(X*BBB(j)-B(j)/B(i)*ppp(i,j));
    end
    s2(i)=0;
    for j=1:n
        s2(i)=s2(i)+alpha(j);
    end
    s2(i)=X*exp(-r*T-s2(i)*deltat);
    s(i)=BB(i)*(s1(i)+s2(i));
end
plot(s)
hold on;
plot(s0);
\end{lstlisting}


\include{body/chap03}
%%
%%
%% 其它部分
%%\backmatter           %%会导致附录的编号A、B、C消失
%% 致谢
%% 搞了好久的致谢等的页眉，没搞定。重定义plain风格，让后续页面出现页眉。dirty code！
\fancypagestyle{plain}{}

%%% Local Variables:
%%% mode: latex
%%% TeX-master: "../main"
%%% End:

\begin{ack}
  衷心感谢导师 xxx 教授对本人的精心指导。他们的言传身教将使
  我终生受益。

  本课题承蒙国家自然科学基金资助，特此致谢。

  感谢 \whutthesis，它的存在让我的论文写作轻松自在了许多，让我的论文格式规整漂亮了
  许多。
\end{ack}

%%
%%
%% 参考文献
\bibliographystyle{GBT7714-2005NLang}
\bibliography{reference/refs}
\addcontentsline{toc}{chapter}{参考文献}
%%
%%
%% 个人简历
\begin{resume}

  \resumeitem{个人简历}

  xxxx 年 xx 月 xx 日出生于 xx 省 xx 县。
  
  xxxx 年 9 月考入 xx 大学 xx 系 xx 专业，xxxx 年 7 月本科毕业并获得 xx 学士学位。
  
  xxxx 年 9 月免试进入 xx 大学 xx 系攻读 xx 学位至今。

  \resumeitem{发表的学术论文} % 发表的和录用的合在一起

  \begin{enumerate}[{[}1{]}]
  \item Yang Y, Ren T L, Zhang L T, et al. Miniature microphone with silicon-
    based ferroelectric thin films. Integrated Ferroelectrics, 2003,
    52:229-235. (SCI 收录, 检索号:758FZ.)
  \item 杨轶, 张宁欣, 任天令, 等. 硅基铁电微声学器件中薄膜残余应力的研究. 中国机
    械工程, 2005, 16(14):1289-1291. (EI 收录, 检索号:0534931 2907.)
  \item 杨轶, 张宁欣, 任天令, 等. 集成铁电器件中的关键工艺研究. 仪器仪表学报,
    2003, 24(S4):192-193. (EI 源刊.)
  \item Yang Y, Ren T L, Zhu Y P, et al. PMUTs for handwriting recognition. In
    press. (已被 Integrated Ferroelectrics 录用. SCI 源刊.)
  \item Wu X M, Yang Y, Cai J, et al. Measurements of ferroelectric MEMS
    microphones. Integrated Ferroelectrics, 2005, 69:417-429. (SCI 收录, 检索号
    :896KM.)
  \item 贾泽, 杨轶, 陈兢, 等. 用于压电和电容微麦克风的体硅腐蚀相关研究. 压电与声
    光, 2006, 28(1):117-119. (EI 收录, 检索号:06129773469.)
  \item 伍晓明, 杨轶, 张宁欣, 等. 基于MEMS技术的集成铁电硅微麦克风. 中国集成电路, 
    2003, 53:59-61.
  \end{enumerate}

  \resumeitem{研究成果} % 有就写，没有就删除
  \begin{enumerate}[{[}1{]}]
  \item 任天令, 杨轶, 朱一平, 等. 硅基铁电微声学传感器畴极化区域控制和电极连接的
    方法: 中国, CN1602118A. (中国专利公开号.)
  \item Ren T L, Yang Y, Zhu Y P, et al. Piezoelectric micro acoustic sensor
    based on ferroelectric materials: USA, No.11/215, 102. (美国发明专利申请号.)
  \end{enumerate}
\end{resume}

%%
%%
%% 几个索引
\makeatletter
%%   插图索引
%%  \listoffigures
%%   表格索引
%%  \listoftables
%%   公式索引
%%  \listofequations
\makeatother
%%
%%
%% 附录
\begin{appendix}
    %%% Local Variables: 
%%% mode: latex
%%% TeX-master: "../main"
%%% End: 

\chapter{外文资料原文}
\label{cha:engorg}
As one of the most widely used techniques in operations research, {\em
  mathematical programming} is defined as a means of maximizing a quantity known
as {\em objective function}, subject to a set of constraints represented by
equations and inequalities. Some known subtopics of mathematical programming are
linear programming, nonlinear programming, multiobjective programming, goal
programming, dynamic programming, and multilevel programming$^{[1]}$.

It is impossible to cover in a single chapter every concept of mathematical
programming. This chapter introduces only the basic concepts and techniques of
mathematical programming such that readers gain an understanding of them
throughout the book$^{[2,3]}$.


\section{Single-Objective Programming}
The general form of single-objective programming (SOP) is written
as follows,
\begin{equation}\tag*{(123)} % 如果附录中的公式不想让它出现在公式索引中，那就请
                             % 用 \tag*{xxxx}
\left\{\begin{array}{l}
\max \,\,f(x)\\[0.1 cm]
\mbox{subject to:} \\ [0.1 cm]
\qquad g_j(x)\le 0,\quad j=1,2,\cdots,p
\end{array}\right.
\end{equation}
which maximizes a real-valued function $f$ of
$x=(x_1,x_2,\cdots,x_n)$ subject to a set of constraints.

\newtheorem{mpdef}{Definition}[chapter]
\begin{mpdef}
In SOP, we call $x$ a decision vector, and
$x_1,x_2,\cdots,x_n$ decision variables. The function
$f$ is called the objective function. The set
\begin{equation}\tag*{(456)} % 这里同理，其它不再一一指定。
S=\left\{x\in\Re^n\bigm|g_j(x)\le 0,\,j=1,2,\cdots,p\right\}
\end{equation}
is called the feasible set. An element $x$ in $S$ is called a
feasible solution.
\end{mpdef}

\newtheorem{mpdefop}[mpdef]{Definition}
\begin{mpdefop}
A feasible solution $x^*$ is called the optimal
solution of SOP if and only if
\begin{equation}
f(x^*)\ge f(x)
\end{equation}
for any feasible solution $x$.
\end{mpdefop}

One of the outstanding contributions to mathematical programming was known as
the Kuhn-Tucker conditions\ref{eq:ktc}. In order to introduce them, let us give
some definitions. An inequality constraint $g_j(x)\le 0$ is said to be active at
a point $x^*$ if $g_j(x^*)=0$. A point $x^*$ satisfying $g_j(x^*)\le 0$ is said
to be regular if the gradient vectors $\nabla g_j(x)$ of all active constraints
are linearly independent.

Let $x^*$ be a regular point of the constraints of SOP and assume that all the
functions $f(x)$ and $g_j(x),j=1,2,\cdots,p$ are differentiable. If $x^*$ is a
local optimal solution, then there exist Lagrange multipliers
$\lambda_j,j=1,2,\cdots,p$ such that the following Kuhn-Tucker conditions hold,
\begin{equation}
\label{eq:ktc}
\left\{\begin{array}{l}
    \nabla f(x^*)-\sum\limits_{j=1}^p\lambda_j\nabla g_j(x^*)=0\\[0.3cm]
    \lambda_jg_j(x^*)=0,\quad j=1,2,\cdots,p\\[0.2cm]
    \lambda_j\ge 0,\quad j=1,2,\cdots,p.
\end{array}\right.
\end{equation}
If all the functions $f(x)$ and $g_j(x),j=1,2,\cdots,p$ are convex and
differentiable, and the point $x^*$ satisfies the Kuhn-Tucker conditions
(\ref{eq:ktc}), then it has been proved that the point $x^*$ is a global optimal
solution of SOP.

\subsection{Linear Programming} 
\label{sec:lp}

If the functions $f(x),g_j(x),j=1,2,\cdots,p$ are all linear, then SOP is called
a {\em linear programming}.

The feasible set of linear is always convex. A point $x$ is called an extreme
point of convex set $S$ if $x\in S$ and $x$ cannot be expressed as a convex
combination of two points in $S$. It has been shown that the optimal solution to
linear programming corresponds to an extreme point of its feasible set provided
that the feasible set $S$ is bounded. This fact is the basis of the {\em simplex
  algorithm} which was developed by Dantzig as a very efficient method for
solving linear programming.
\begin{table}[ht]
\centering
  \centering
  \caption*{Table~1\hskip1em This is an example for manually numbered table, which
    would not appear in the list of tables}
  \label{tab:badtabular2}
  \begin{tabular}[c]{|c|m{0.8in}|c|c|c|c|c|}\hline
    \multicolumn{2}{|c|}{Network Topology} & \# of nodes & 
    \multicolumn{3}{c|}{\# of clients} & Server \\\hline
    GT-ITM & Waxman Transit-Stub & 600 &
    \multirow{2}{2em}{2\%}& 
    \multirow{2}{2em}{10\%}& 
    \multirow{2}{2em}{50\%}& 
    \multirow{2}{1.2in}{Max. Connectivity}\\\cline{1-3}
    \multicolumn{2}{|c|}{Inet-2.1} & 6000 & & & &\\\hline
    \multirow{2}{1in}{Xue} & Rui  & Ni &\multicolumn{4}{c|}{\multirow{2}*{\whutthesis}}\\\cline{2-3}
    & \multicolumn{2}{c|}{ABCDEF} &\multicolumn{4}{c|}{} \\\hline
\end{tabular}  
\end{table}

Roughly speaking, the simplex algorithm examines only the extreme points of the
feasible set, rather than all feasible points. At first, the simplex algorithm
selects an extreme point as the initial point. The successive extreme point is
selected so as to improve the objective function value. The procedure is
repeated until no improvement in objective function value can be made. The last
extreme point is the optimal solution.

\subsection{Nonlinear Programming}

If at least one of the functions $f(x),g_j(x),j=1,2,\cdots,p$ is nonlinear, then
SOP is called a {\em nonlinear programming}.

A large number of classical optimization methods have been developed to treat
special-structural nonlinear programming based on the mathematical theory
concerned with analyzing the structure of problems.
\begin{figure}[h]
  \centering
  \includegraphics[clip]{thu-lib-logo}
  \caption*{Figure~1\hskip1em This is an example for manually numbered figure,
    which would not appear in the list of figures}
  \label{tab:badfigure2}    
\end{figure}

Now we consider a nonlinear programming which is confronted solely with
maximizing a real-valued function with domain $\Re^n$.  Whether derivatives are
available or not, the usual strategy is first to select a point in $\Re^n$ which
is thought to be the most likely place where the maximum exists. If there is no
information available on which to base such a selection, a point is chosen at
random. From this first point an attempt is made to construct a sequence of
points, each of which yields an improved objective function value over its
predecessor. The next point to be added to the sequence is chosen by analyzing
the behavior of the function at the previous points. This construction continues
until some termination criterion is met. Methods based upon this strategy are
called {\em ascent methods}, which can be classified as {\em direct methods},
{\em gradient methods}, and {\em Hessian methods} according to the information
about the behavior of objective function $f$. Direct methods require only that
the function can be evaluated at each point. Gradient methods require the
evaluation of first derivatives of $f$. Hessian methods require the evaluation
of second derivatives. In fact, there is no superior method for all
problems. The efficiency of a method is very much dependent upon the objective
function.

\subsection{Integer Programming}

{\em Integer programming} is a special mathematical programming in which all of
the variables are assumed to be only integer values. When there are not only
integer variables but also conventional continuous variables, we call it {\em
  mixed integer programming}. If all the variables are assumed either 0 or 1,
then the problem is termed a {\em zero-one programming}. Although integer
programming can be solved by an {\em exhaustive enumeration} theoretically, it
is impractical to solve realistically sized integer programming problems. The
most successful algorithm so far found to solve integer programming is called
the {\em branch-and-bound enumeration} developed by Balas (1965) and Dakin
(1965). The other technique to integer programming is the {\em cutting plane
  method} developed by Gomory (1959).

\hfill\textit{Uncertain Programming\/}\quad(\textsl{BaoDing Liu, 2006.2})

\section*{References}
\noindent{\itshape NOTE: these references are only for demonstration, they are
  not real citations in the original text.}

\begin{enumerate}[{$[$}1{$]$}]
\item Donald E. Knuth. The \TeX book. Addison-Wesley, 1984. ISBN: 0-201-13448-9
\item Paul W. Abrahams, Karl Berry and Kathryn A. Hargreaves. \TeX\ for the
  Impatient. Addison-Wesley, 1990. ISBN: 0-201-51375-7
\item David Salomon. The advanced \TeX book.  New York : Springer, 1995. ISBN:0-387-94556-3
\end{enumerate}

\chapter{外文资料的调研阅读报告或书面翻译}
\section{单目标规划}
北冥有鱼，其名为鲲。鲲之大，不知其几千里也。化而为鸟，其名为鹏。鹏之背，不知其几
千里也。怒而飞，其翼若垂天之云。是鸟也，海运则将徙于南冥。南冥者，天池也。 
\begin{equation}\tag*{(123)}
 p(y|\mathbf{x}) = \frac{p(\mathbf{x},y)}{p(\mathbf{x})}=
\frac{p(\mathbf{x}|y)p(y)}{p(\mathbf{x})}
\end{equation}

吾生也有涯，而知也无涯。以有涯随无涯，殆已！已而为知者，殆而已矣！为善无近名，为
恶无近刑，缘督以为经，可以保身，可以全生，可以养亲，可以尽年。

\subsection{线性规划}
庖丁为文惠君解牛，手之所触，肩之所倚，足之所履，膝之所倚，砉然响然，奏刀騞然，莫
不中音，合于桑林之舞，乃中经首之会。
\begin{table}[ht]
\centering
  \centering
  \caption*{表~1\hskip1em 这是手动编号但不出现在索引中的一个表格例子}
  \label{tab:badtabular3}
  \begin{tabular}[c]{|c|m{0.8in}|c|c|c|c|c|}\hline
    \multicolumn{2}{|c|}{Network Topology} & \# of nodes & 
    \multicolumn{3}{c|}{\# of clients} & Server \\\hline
    GT-ITM & Waxman Transit-Stub & 600 &
    \multirow{2}{2em}{2\%}& 
    \multirow{2}{2em}{10\%}& 
    \multirow{2}{2em}{50\%}& 
    \multirow{2}{1.2in}{Max. Connectivity}\\\cline{1-3}
    \multicolumn{2}{|c|}{Inet-2.1} & 6000 & & & &\\\hline
    \multirow{2}{1in}{Xue} & Rui  & Ni &\multicolumn{4}{c|}{\multirow{2}*{\whutthesis}}\\\cline{2-3}
    & \multicolumn{2}{c|}{ABCDEF} &\multicolumn{4}{c|}{} \\\hline
\end{tabular}  
\end{table}

文惠君曰：“嘻，善哉！技盖至此乎？”庖丁释刀对曰：“臣之所好者道也，进乎技矣。始臣之
解牛之时，所见无非全牛者；三年之后，未尝见全牛也；方今之时，臣以神遇而不以目视，
官知止而神欲行。依乎天理，批大郤，导大窾，因其固然。技经肯綮之未尝，而况大坬乎！
良庖岁更刀，割也；族庖月更刀，折也；今臣之刀十九年矣，所解数千牛矣，而刀刃若新发
于硎。彼节者有间而刀刃者无厚，以无厚入有间，恢恢乎其于游刃必有余地矣。是以十九年
而刀刃若新发于硎。虽然，每至于族，吾见其难为，怵然为戒，视为止，行为迟，动刀甚微，
謋然已解，如土委地。提刀而立，为之而四顾，为之踌躇满志，善刀而藏之。”

文惠君曰：“善哉！吾闻庖丁之言，得养生焉。”


\subsection{非线性规划}
孔子与柳下季为友，柳下季之弟名曰盗跖。盗跖从卒九千人，横行天下，侵暴诸侯。穴室枢
户，驱人牛马，取人妇女。贪得忘亲，不顾父母兄弟，不祭先祖。所过之邑，大国守城，小
国入保，万民苦之。孔子谓柳下季曰：“夫为人父者，必能诏其子；为人兄者，必能教其弟。
若父不能诏其子，兄不能教其弟，则无贵父子兄弟之亲矣。今先生，世之才士也，弟为盗
跖，为天下害，而弗能教也，丘窃为先生羞之。丘请为先生往说之。”
\begin{figure}[h]
  \centering
  \includegraphics{hello}
  \caption*{图~1\hskip1em 这是手动编号但不出现索引中的图片的例子}
  \label{tab:badfigure3}    
\end{figure}

柳下季曰：“先生言为人父者必能诏其子，为人兄者必能教其弟，若子不听父之诏，弟不受
兄之教，虽今先生之辩，将奈之何哉？且跖之为人也，心如涌泉，意如飘风，强足以距敌，
辩足以饰非。顺其心则喜，逆其心则怒，易辱人以言。先生必无往。”

孔子不听，颜回为驭，子贡为右，往见盗跖。

\subsection{整数规划}
盗跖乃方休卒徒大山之阳，脍人肝而餔之。孔子下车而前，见谒者曰：“鲁人孔丘，闻将军
高义，敬再拜谒者。”谒者入通。盗跖闻之大怒，目如明星，发上指冠，曰：“此夫鲁国之
巧伪人孔丘非邪？为我告之：尔作言造语，妄称文、武，冠枝木之冠，带死牛之胁，多辞缪
说，不耕而食，不织而衣，摇唇鼓舌，擅生是非，以迷天下之主，使天下学士不反其本，妄
作孝弟，而侥幸于封侯富贵者也。子之罪大极重，疾走归！不然，我将以子肝益昼餔之膳。”


\chapter{其它附录}
其它附录的内容可以放到这里，当然如果你愿意，可
以把这部分也放到独立的文件中，然后将其 \verb|\input| 到主文件中。
\end{appendix}
%%
%%
\end{document}
%%
%% END of mythsis.tex
%</myth>
%    \end{macrocode}
% \end{example}
%
% \subsection{选项}
% \label{sec:option}
% 本模板提供了一些选项以方便使用：
% \begin{description}
%
% \item[master]
%   如果写硕士论文将此选项打开。
%   \begin{example}
% \documentclass[master]{WHUTthesis}
%   \end{example}
%
% \item[doctor]
%   如果写博士论文将此选项打开。
%   \begin{example}
% \documentclass[doctor]{WHUTthesis}
%   \end{example}
%
% \item[openany, openright]
%   正规出版物的章节出现在奇数页，也就是右手边的页面，这就是 \texttt{openright}。
%   在这种情况下，如果前一章的最后一页也是奇数，那么
%   模板会自动生成一个纯粹的空白页，很多人不是很习惯这种方式，而且学校的格式似乎
%   更倾向于页面连续，那就是通常所说的 \texttt{openany}。{\fangsong 目前所有论文都是
%      openany。}这两个选项不用专门设置，\whutthesis{} 会根据当前论文类型自动选择。
%   \begin{example}
% \documentclass[doctor,twoside,openright]{WHUTthesis}
%   \end{example}
%
% \item[winfonts, adobefonts]
%   winfonts选项使得在Windows平台下使用系统提供的字体。另外Adobe公司提供了四款免费
%   的OTF字体（宋体、仿宋、黑体、楷体），一般来说OTF字体的显示效果要优于ttf字体，使用
%   adobefonts可以打开OTF字体。
%   \begin{example}
% \documentclass[master,oneside,adobefonts]{WHUTthesis}
%   \end{example}
%
% \item[composing]
%   论文在写作时和最终提交的版本可能存在一些差别，如写作时给自己看的一些批注。
%   提供此选项以便打开写作模式。之所以不用draft选项是因为该选项下不会加载图片和代码列表
%   等内容，不便于观察排版结果。\textbf{注意：由composing选项切换到无该选项编译文档时可能
%   会出现pgf宏包的错误，这可能是因为上次带composing选项编译时中间文件中存在使用mynotes
%   命令而由\textsf{todonotes}宏包引入的pgf指令，简单跳过编译或者清除上次编译的中间文件然后执行
%   全新的编译即能顺利进行。}
%   \begin{example}
% \documentclass[master,oneside,composing]{WHUTthesis}
%   \end{example}
% \end{description}
%
% \subsection{命令}
% \label{sec:command}
% 模板中的命令分为两类：一是格式控制，二是内容替换。格式控制如字体、字号、字距和
% 行距。内容替换如姓名、院系、专业、致谢等等。其中内容替换命令居多，而且主要集中
% 在封面上。首先来看格式控制命令。
%
% \subsubsection{基本控制命令}
% \label{sec:basiccom}
%
% \myentry{字体}字体采用ctex的默认字体选择命令
% \DescribeMacro{\songti}
% \DescribeMacro{\fangsong}
% \DescribeMacro{\heiti}
% \DescribeMacro{\kaishu}
% \DescribeMacro{\lishu}
% \DescribeMacro{\youyuan}
% 等分别用来切换宋体、仿宋、黑体、楷体、隶书和幼圆字体。
%
% \begin{example}
% {\songti 乾：元，亨，利贞}
% {\fangsong 初九，潜龙勿用}
% {\heiti 九二，见龙在田，利见大人}
% {\kaishu 九三，君子终日乾乾，夕惕若，厉，无咎}
% {\lishu 九四，或跃在渊，无咎}
% \end{example}
%
% \myentry{字号}
% \DescribeMacro{\zihao}
% 命令为ctex中用来选择字号的命令。其中参数n为要使用的字号，
% 在n前加负号-表示小号字体，
% 使用方法为：\cs{command}\oarg{num}。学校论文要求的小四号在模板中已经设定好了。
%
% \begin{example}
% {\zihao{2} 二号 \zihao{3} 三号 \zihao{4} 四号 \zihao{-4} 小四 \zihao{7} 七号}
% \end{example}
%
% \myentry{字距}
% \DescribeMacro{\ziju}
% 更改汉字之间默认的距离，使用格式为 |\ziju{4bp}|，其中的距离只要是合格的 \TeX{} 距离即可。
%
% \myentry{密级}
% \DescribeMacro{\secretlevel}
% 定义秘密级别：
%   \begin{example}
% \secretlevel{内部}
%   \end{example}
%
% \myentry{引用方式}
% \DescribeMacro{\onlinecite}
%
% 学校要求的参考文献引用有两种模式：（1）上标模式。比如``同样的工作有很
% 多$^{[1,2]}$\ldots''。（2）正文模式。比如``文[3] 中详细说明了\ldots''。其中上标
% 模式使用远比正文模式频繁，所以为了符合使用习惯，上标模式仍然用常规
% 的 |\cite{key}|，而 |\onlinecite{key}| 则用来生成正文模式。
%
% 关于参考文献模板推荐使用 BIB\TeX，关于中文参考文献需要额外增加一个 Entry: lang，将其设置为 \texttt{zh} 
% 用来指示此参考文献为中文，以便 thubib.bst 处理。如：
% \begin{example}
% @INPROCEEDINGS{cnproceed,
%   author    = {王重阳 and 黄药师 and 欧阳峰 and 洪七公 and 段皇帝},
%   title     = {武林高手从入门到精通},
%   booktitle = {第~$N$~次华山论剑},
%   year      = 2006,
%   address   = {西安, 中国},
%   month     = sep,
%   lang      = "zh",
% }
%
% @ARTICLE{cnarticle,
%   AUTHOR  = "贾宝玉 and 林黛玉 and 薛宝钗 and 贾探春",
%   TITLE   = "论刘姥姥食量大如牛之现实意义",
%   JOURNAL = "红楼梦杂谈",
%   PAGES   = "260--266",
%   VOLUME  = "224",
%   YEAR    = "1800",
%   LANG    = "zh",
% }
% \end{example}
%
%
% \myentry{破折号} 
% \DescribeMacro{\pozhehao}
% 中文破折号在 CJK-\LaTeX\ 里没有很好的处理，我们平时输入的都是两个小短线，比如这
% 样，{\heiti 中国——中华人民共和国}。这不符合中文习惯。所以这里定义了一个命令生成更
% 好看的破折号，不过这似乎不是一个好的解决办法。有同学说不能用在 |\section| 等命
% 令中使用，简单的办法是可以提供一个不带破折号的段标题：\cs{section}\oarg{没有破
%   折号精简标题}\marg{带破折号的标题}。
%
%
% \subsubsection{封面命令}
% \label{sec:titlepage}
% 下面是内容替换命令，其中以 |c| 开头的命令跟中文相关，|e| 开头则为对应的英文。
% 这部分的命令数目比较多，但实际上都相当简单，套用即可。
%
% 大多数命令的使用方法都是： \cs{command}\marg{arg}，例外者将具体指出。这些命令都
% 在示例文档的 data/cover.tex 中。
%
% \myentry{论文标题}
% \DescribeMacro{\ctitle}
% \DescribeMacro{\etitle}
% \DescribeMacro{\ctitleformated}
% \DescribeMacro{\etitleformated}
% \begin{example}
% \ctitle{论文中文题目}
% \etitle{Thesis English Title}
% \ctitleformated{论文题目太长了\\ 需要换行}
% \etitleformated{The title is so long that\\ it must be devidec into two lines}
% \end{example}
%
% \myentry{作者姓名}
% \DescribeMacro{\cauthor}
% \DescribeMacro{\eauthor}
% \begin{example}
% \cauthor{中文姓名}
% \eauthor{Your name in PinYin}
% \end{example}
%
% \myentry{申请学位名称}
% \DescribeMacro{\cdegree}
% \DescribeMacro{\edegree}
% \begin{example}
% \cdegree{您要申请什么学位}
% \edegree{degree in English}
% \end{example}
%
% \myentry{院系名称}
% \DescribeMacro{\cauthordepartment}
% \DescribeMacro{\eauthordepartment}
%
% \cs{cauthordepartment} 可以加一个可选参数，如：\cs{cauthordepartmentl}\oarg{精简}\marg{详
%   细}，在填写空间有限时可以采用精简名称。如\textbf{计算机科学与技术}和\textbf{计算机}。
% \begin{example}
% \cauthordepartment[系名简称]{系名全称}
% \eauthordepartment{Department}
% \end{example}
%
% \myentry{专业名称}
% \DescribeMacro{\cmajor}
% \DescribeMacro{\emajor}
% \begin{example}
% \cmajor{专业名称}
% \emajor{Major in English}
% \end{example}
%
% \myentry{导师姓名}
% \DescribeMacro{\cadvisorname}
% \DescribeMacro{\eadvisorname}
% \begin{example}
% \cadvisorname{大老板}
% \eadvisorname{BIG BOSS}
% \end{example}
%
% \myentry{导师职称}
% \DescribeMacro{\cadvisortitle}
% \begin{example}
% \cadvisortitle{教授}
% \end{example}
%
% \myentry{导师学位}
% \DescribeMacro{\cadvisordegree}
% \begin{example}
% \cadvisordegree{博士}
% \end{example}
%
% \myentry{导师单位}
% \DescribeMacro{\cadvisordepartment}
% \begin{example}
% \cadvisordepartment{计算机科学与技术学院}
% \end{example}
%
% \myentry{导师邮编}
% \DescribeMacro{\advisorzipcode}
% \begin{example}
% \advisorzipcode{430063}
% \end{example}
%
% \myentry{论文提交日期}
% \DescribeMacro{\csubmitdate}
% \begin{example}
% \csubmitdate{\the\year 年\the\month 月}
% \end{example}
%
% \myentry{论文答辩日期}
% \DescribeMacro{\coraldefensedate}
% \begin{example}
% \coraldefensedate{\the\year 年\the\month 月}
% \end{example}
%
% \myentry{学位授予单位}
% \DescribeMacro{\cauthorizedepartment}
% \begin{example}
% \cauthorizedepartment{武汉理工大学}
% \end{example}
%
% \myentry{学位授予日期}
% \DescribeMacro{\cauthorizedate}
% \begin{example}
% \cauthorizedate{\the\year 年\the\month 月}
% \end{example}
%
% \myentry{答辩委员会主席}
% \DescribeMacro{\cchairman}
% \begin{example}
% \cchairman{胡卫谊}
% \end{example}
%
% \myentry{评阅人}
% \DescribeMacro{\cannotatora}
% \DescribeMacro{\cannotatorb}
% \begin{example}
% \cannotatora{胡卫谊a}
% \cannotatorb{胡卫谊b}
% \end{example}
%
% \myentry{副导师姓名}
% \DescribeMacro{\ccoadvisorname}
% \DescribeMacro{\ecoadvisorname}
% 副指导教师。
% \begin{example}
% \ccoadvisorname{小老板}
% \ecoadvisorname{small boss}
% \end{example}
%
% \myentry{论文成文日期}
% \DescribeMacro{\cdate}
% \DescribeMacro{\edate}
% 默认为当前时间，也可以自己指定。
% \begin{example}
% \cdate{中文日期}
% \edate{English Date}
% \end{example}
%
% \myentry{摘要}
% \DescribeEnv{cabstract}
% \DescribeEnv{eabstract}
% \begin{example}
% \begin{cabstract}
%  摘要请写在这里...
% \end{cabstract}
% \begin{eabstract}
%  here comes English abstract...
% \end{eabstract}
% \end{example}
%
% \myentry{关键词}
% \DescribeMacro{\ckeywords}
% \DescribeMacro{\ekeywords}
% 关键词用英文逗号分割写入相应的命令中，模板会解析各关键词并生成符合不同论文格式
% 要求的关键词格式。
% \begin{example}
% \ckeywords{关键词 1, 关键词 2}
% \ekeywords{keyword 1, key word 2}
% \end{example}
%
% \subsubsection{其它部分}
% \label{sec:otherparts}
% 论文其它主要部分命令：
%
% \myentry{符号对照表}
% \DescribeEnv{denotation}
% 主要符号表环境。简单定义的一个 list，跟 description 非常类似，使用方法参见示例
% 文件。带一个可选参数，用来指定符号列的宽度（默认为 2.5cm）。
% \begin{example}
% \begin{denotation}
%   \item[E] 能量
%   \item[m] 质量
%   \item[c] 光速
% \end{denotation}
% \end{example}
%
% 如果你觉得符号列的宽度不满意，那可以这样来调整：
% \begin{example}
% \begin{denotation}[1.5cm] % 设置为 1.5cm
%   \item[E] 能量
%   \item[m] 质量
%   \item[c] 光速
% \end{denotation}
% \end{example}
%
% \myentry{索引}
% 插图、表格和公式三个索引命令分别如下，将其插入到期望的位置即可（带星号的命令表
% 示对应的索引表不会出现在目录中）：
%
% \begin{center}
% \begin{tabular}{ll}
% \hline
%   {\heiti 命令} & {\heiti 说明} \\\hline
% \cs{listoffigures} & 插图索引\\
% \cs{listoffigures*} & \\\hline
% \cs{listoftables} & 表格索引\\
% \cs{listoftables*} & \\\hline
% \cs{listofequations} & 公式索引\\
% \cs{listofequations*} & \\\hline
% \end{tabular}
% \end{center}
%
% \LaTeX{} 默认支持插图和表格索引，是通过 \cs{caption} 命令完成的，因此它们必须出
% 现在浮动环境中，否则不被计数。
%
% 有的同学不想让某个表格或者图片出现在索引里面，那么请使用命令 \cs{caption*}，这
% 个命令不会给表格编号，也就是出来的只有标题文字而没有``表~xx''，``图~xx''，否则
% 索引里面序号不连续就显得不伦不类，这也是 \LaTeX{} 里星号命令默认的规则。
%
% 公式索引为本模板扩展，模板扩展了 \pkg{amsmath} 几个内部命令，使得公式编号样式和
% 自动索引功能非常方便。一般来说，你用到的所有数学环境编号都没问题了，这个可以参
% 看示例文档。如果你有个非常特殊的数学环境需要加入公式索引，那么请使
% 用 \cs{equcaption}\marg{编号}。此命令表示 equation caption，带一个参数，即显示
% 在索引中的编号。因为公式与图表不同，我们很少给一个公式附加一个标题，之所以起这
% 么个名字是因为图表就是通过 \cs{caption} 加入索引的，\cs{equcaption} 完全就是为
% 了生成公式列表，不产生什么标题。
%
% 使用方法如下。假如有一个非 equation 数学环境 mymath，只要在其中写一
% 句 \cs{equcaption} 就可以将它加入公式列表。
% \begin{example}
% \begin{mymath}
%   \label{eq:emc2}\equcaption{\ref{eq:emc2}}
%   E=mc^2
% \end{mymath}
% \end{example}
%
% 当然 mymath 正文中公式的编号需要你自己来做。
%
% 同图表一样，附录中的公式有时候也不希望它跟全文统一编号，而且不希望它出现在公式
% 索引中，目前的解决办法就是利用 \cs{tag*}\marg{公式编号} 来解决。用法很简单，此
% 处不再罗嗦，实例请参看示例文档附录 A 的前两个公式。
%
% \myentry{简历}
% \DescribeEnv{resume}\DescribeMacro{\resumeitem}
% 开启个人简历章节，包括发表文章列表等。其实就是一个 chapter。里面的每个子项目请用命令 |\resumeitem{sub title}|。
%
% 这里就不再列举例子了，请参看示例文档的 body/resume.tex。
%
% \myentry{附录}
% \DescribeEnv{appendix}
% 所有的附录都插到这里来。
% \begin{example}
% \begin{appendix}
%  %%% Local Variables: 
%%% mode: latex
%%% TeX-master: "../main"
%%% End: 

\chapter{外文资料原文}
\label{cha:engorg}
As one of the most widely used techniques in operations research, {\em
  mathematical programming} is defined as a means of maximizing a quantity known
as {\em objective function}, subject to a set of constraints represented by
equations and inequalities. Some known subtopics of mathematical programming are
linear programming, nonlinear programming, multiobjective programming, goal
programming, dynamic programming, and multilevel programming$^{[1]}$.

It is impossible to cover in a single chapter every concept of mathematical
programming. This chapter introduces only the basic concepts and techniques of
mathematical programming such that readers gain an understanding of them
throughout the book$^{[2,3]}$.


\section{Single-Objective Programming}
The general form of single-objective programming (SOP) is written
as follows,
\begin{equation}\tag*{(123)} % 如果附录中的公式不想让它出现在公式索引中，那就请
                             % 用 \tag*{xxxx}
\left\{\begin{array}{l}
\max \,\,f(x)\\[0.1 cm]
\mbox{subject to:} \\ [0.1 cm]
\qquad g_j(x)\le 0,\quad j=1,2,\cdots,p
\end{array}\right.
\end{equation}
which maximizes a real-valued function $f$ of
$x=(x_1,x_2,\cdots,x_n)$ subject to a set of constraints.

\newtheorem{mpdef}{Definition}[chapter]
\begin{mpdef}
In SOP, we call $x$ a decision vector, and
$x_1,x_2,\cdots,x_n$ decision variables. The function
$f$ is called the objective function. The set
\begin{equation}\tag*{(456)} % 这里同理，其它不再一一指定。
S=\left\{x\in\Re^n\bigm|g_j(x)\le 0,\,j=1,2,\cdots,p\right\}
\end{equation}
is called the feasible set. An element $x$ in $S$ is called a
feasible solution.
\end{mpdef}

\newtheorem{mpdefop}[mpdef]{Definition}
\begin{mpdefop}
A feasible solution $x^*$ is called the optimal
solution of SOP if and only if
\begin{equation}
f(x^*)\ge f(x)
\end{equation}
for any feasible solution $x$.
\end{mpdefop}

One of the outstanding contributions to mathematical programming was known as
the Kuhn-Tucker conditions\ref{eq:ktc}. In order to introduce them, let us give
some definitions. An inequality constraint $g_j(x)\le 0$ is said to be active at
a point $x^*$ if $g_j(x^*)=0$. A point $x^*$ satisfying $g_j(x^*)\le 0$ is said
to be regular if the gradient vectors $\nabla g_j(x)$ of all active constraints
are linearly independent.

Let $x^*$ be a regular point of the constraints of SOP and assume that all the
functions $f(x)$ and $g_j(x),j=1,2,\cdots,p$ are differentiable. If $x^*$ is a
local optimal solution, then there exist Lagrange multipliers
$\lambda_j,j=1,2,\cdots,p$ such that the following Kuhn-Tucker conditions hold,
\begin{equation}
\label{eq:ktc}
\left\{\begin{array}{l}
    \nabla f(x^*)-\sum\limits_{j=1}^p\lambda_j\nabla g_j(x^*)=0\\[0.3cm]
    \lambda_jg_j(x^*)=0,\quad j=1,2,\cdots,p\\[0.2cm]
    \lambda_j\ge 0,\quad j=1,2,\cdots,p.
\end{array}\right.
\end{equation}
If all the functions $f(x)$ and $g_j(x),j=1,2,\cdots,p$ are convex and
differentiable, and the point $x^*$ satisfies the Kuhn-Tucker conditions
(\ref{eq:ktc}), then it has been proved that the point $x^*$ is a global optimal
solution of SOP.

\subsection{Linear Programming} 
\label{sec:lp}

If the functions $f(x),g_j(x),j=1,2,\cdots,p$ are all linear, then SOP is called
a {\em linear programming}.

The feasible set of linear is always convex. A point $x$ is called an extreme
point of convex set $S$ if $x\in S$ and $x$ cannot be expressed as a convex
combination of two points in $S$. It has been shown that the optimal solution to
linear programming corresponds to an extreme point of its feasible set provided
that the feasible set $S$ is bounded. This fact is the basis of the {\em simplex
  algorithm} which was developed by Dantzig as a very efficient method for
solving linear programming.
\begin{table}[ht]
\centering
  \centering
  \caption*{Table~1\hskip1em This is an example for manually numbered table, which
    would not appear in the list of tables}
  \label{tab:badtabular2}
  \begin{tabular}[c]{|c|m{0.8in}|c|c|c|c|c|}\hline
    \multicolumn{2}{|c|}{Network Topology} & \# of nodes & 
    \multicolumn{3}{c|}{\# of clients} & Server \\\hline
    GT-ITM & Waxman Transit-Stub & 600 &
    \multirow{2}{2em}{2\%}& 
    \multirow{2}{2em}{10\%}& 
    \multirow{2}{2em}{50\%}& 
    \multirow{2}{1.2in}{Max. Connectivity}\\\cline{1-3}
    \multicolumn{2}{|c|}{Inet-2.1} & 6000 & & & &\\\hline
    \multirow{2}{1in}{Xue} & Rui  & Ni &\multicolumn{4}{c|}{\multirow{2}*{\whutthesis}}\\\cline{2-3}
    & \multicolumn{2}{c|}{ABCDEF} &\multicolumn{4}{c|}{} \\\hline
\end{tabular}  
\end{table}

Roughly speaking, the simplex algorithm examines only the extreme points of the
feasible set, rather than all feasible points. At first, the simplex algorithm
selects an extreme point as the initial point. The successive extreme point is
selected so as to improve the objective function value. The procedure is
repeated until no improvement in objective function value can be made. The last
extreme point is the optimal solution.

\subsection{Nonlinear Programming}

If at least one of the functions $f(x),g_j(x),j=1,2,\cdots,p$ is nonlinear, then
SOP is called a {\em nonlinear programming}.

A large number of classical optimization methods have been developed to treat
special-structural nonlinear programming based on the mathematical theory
concerned with analyzing the structure of problems.
\begin{figure}[h]
  \centering
  \includegraphics[clip]{thu-lib-logo}
  \caption*{Figure~1\hskip1em This is an example for manually numbered figure,
    which would not appear in the list of figures}
  \label{tab:badfigure2}    
\end{figure}

Now we consider a nonlinear programming which is confronted solely with
maximizing a real-valued function with domain $\Re^n$.  Whether derivatives are
available or not, the usual strategy is first to select a point in $\Re^n$ which
is thought to be the most likely place where the maximum exists. If there is no
information available on which to base such a selection, a point is chosen at
random. From this first point an attempt is made to construct a sequence of
points, each of which yields an improved objective function value over its
predecessor. The next point to be added to the sequence is chosen by analyzing
the behavior of the function at the previous points. This construction continues
until some termination criterion is met. Methods based upon this strategy are
called {\em ascent methods}, which can be classified as {\em direct methods},
{\em gradient methods}, and {\em Hessian methods} according to the information
about the behavior of objective function $f$. Direct methods require only that
the function can be evaluated at each point. Gradient methods require the
evaluation of first derivatives of $f$. Hessian methods require the evaluation
of second derivatives. In fact, there is no superior method for all
problems. The efficiency of a method is very much dependent upon the objective
function.

\subsection{Integer Programming}

{\em Integer programming} is a special mathematical programming in which all of
the variables are assumed to be only integer values. When there are not only
integer variables but also conventional continuous variables, we call it {\em
  mixed integer programming}. If all the variables are assumed either 0 or 1,
then the problem is termed a {\em zero-one programming}. Although integer
programming can be solved by an {\em exhaustive enumeration} theoretically, it
is impractical to solve realistically sized integer programming problems. The
most successful algorithm so far found to solve integer programming is called
the {\em branch-and-bound enumeration} developed by Balas (1965) and Dakin
(1965). The other technique to integer programming is the {\em cutting plane
  method} developed by Gomory (1959).

\hfill\textit{Uncertain Programming\/}\quad(\textsl{BaoDing Liu, 2006.2})

\section*{References}
\noindent{\itshape NOTE: these references are only for demonstration, they are
  not real citations in the original text.}

\begin{enumerate}[{$[$}1{$]$}]
\item Donald E. Knuth. The \TeX book. Addison-Wesley, 1984. ISBN: 0-201-13448-9
\item Paul W. Abrahams, Karl Berry and Kathryn A. Hargreaves. \TeX\ for the
  Impatient. Addison-Wesley, 1990. ISBN: 0-201-51375-7
\item David Salomon. The advanced \TeX book.  New York : Springer, 1995. ISBN:0-387-94556-3
\end{enumerate}

\chapter{外文资料的调研阅读报告或书面翻译}
\section{单目标规划}
北冥有鱼，其名为鲲。鲲之大，不知其几千里也。化而为鸟，其名为鹏。鹏之背，不知其几
千里也。怒而飞，其翼若垂天之云。是鸟也，海运则将徙于南冥。南冥者，天池也。 
\begin{equation}\tag*{(123)}
 p(y|\mathbf{x}) = \frac{p(\mathbf{x},y)}{p(\mathbf{x})}=
\frac{p(\mathbf{x}|y)p(y)}{p(\mathbf{x})}
\end{equation}

吾生也有涯，而知也无涯。以有涯随无涯，殆已！已而为知者，殆而已矣！为善无近名，为
恶无近刑，缘督以为经，可以保身，可以全生，可以养亲，可以尽年。

\subsection{线性规划}
庖丁为文惠君解牛，手之所触，肩之所倚，足之所履，膝之所倚，砉然响然，奏刀騞然，莫
不中音，合于桑林之舞，乃中经首之会。
\begin{table}[ht]
\centering
  \centering
  \caption*{表~1\hskip1em 这是手动编号但不出现在索引中的一个表格例子}
  \label{tab:badtabular3}
  \begin{tabular}[c]{|c|m{0.8in}|c|c|c|c|c|}\hline
    \multicolumn{2}{|c|}{Network Topology} & \# of nodes & 
    \multicolumn{3}{c|}{\# of clients} & Server \\\hline
    GT-ITM & Waxman Transit-Stub & 600 &
    \multirow{2}{2em}{2\%}& 
    \multirow{2}{2em}{10\%}& 
    \multirow{2}{2em}{50\%}& 
    \multirow{2}{1.2in}{Max. Connectivity}\\\cline{1-3}
    \multicolumn{2}{|c|}{Inet-2.1} & 6000 & & & &\\\hline
    \multirow{2}{1in}{Xue} & Rui  & Ni &\multicolumn{4}{c|}{\multirow{2}*{\whutthesis}}\\\cline{2-3}
    & \multicolumn{2}{c|}{ABCDEF} &\multicolumn{4}{c|}{} \\\hline
\end{tabular}  
\end{table}

文惠君曰：“嘻，善哉！技盖至此乎？”庖丁释刀对曰：“臣之所好者道也，进乎技矣。始臣之
解牛之时，所见无非全牛者；三年之后，未尝见全牛也；方今之时，臣以神遇而不以目视，
官知止而神欲行。依乎天理，批大郤，导大窾，因其固然。技经肯綮之未尝，而况大坬乎！
良庖岁更刀，割也；族庖月更刀，折也；今臣之刀十九年矣，所解数千牛矣，而刀刃若新发
于硎。彼节者有间而刀刃者无厚，以无厚入有间，恢恢乎其于游刃必有余地矣。是以十九年
而刀刃若新发于硎。虽然，每至于族，吾见其难为，怵然为戒，视为止，行为迟，动刀甚微，
謋然已解，如土委地。提刀而立，为之而四顾，为之踌躇满志，善刀而藏之。”

文惠君曰：“善哉！吾闻庖丁之言，得养生焉。”


\subsection{非线性规划}
孔子与柳下季为友，柳下季之弟名曰盗跖。盗跖从卒九千人，横行天下，侵暴诸侯。穴室枢
户，驱人牛马，取人妇女。贪得忘亲，不顾父母兄弟，不祭先祖。所过之邑，大国守城，小
国入保，万民苦之。孔子谓柳下季曰：“夫为人父者，必能诏其子；为人兄者，必能教其弟。
若父不能诏其子，兄不能教其弟，则无贵父子兄弟之亲矣。今先生，世之才士也，弟为盗
跖，为天下害，而弗能教也，丘窃为先生羞之。丘请为先生往说之。”
\begin{figure}[h]
  \centering
  \includegraphics{hello}
  \caption*{图~1\hskip1em 这是手动编号但不出现索引中的图片的例子}
  \label{tab:badfigure3}    
\end{figure}

柳下季曰：“先生言为人父者必能诏其子，为人兄者必能教其弟，若子不听父之诏，弟不受
兄之教，虽今先生之辩，将奈之何哉？且跖之为人也，心如涌泉，意如飘风，强足以距敌，
辩足以饰非。顺其心则喜，逆其心则怒，易辱人以言。先生必无往。”

孔子不听，颜回为驭，子贡为右，往见盗跖。

\subsection{整数规划}
盗跖乃方休卒徒大山之阳，脍人肝而餔之。孔子下车而前，见谒者曰：“鲁人孔丘，闻将军
高义，敬再拜谒者。”谒者入通。盗跖闻之大怒，目如明星，发上指冠，曰：“此夫鲁国之
巧伪人孔丘非邪？为我告之：尔作言造语，妄称文、武，冠枝木之冠，带死牛之胁，多辞缪
说，不耕而食，不织而衣，摇唇鼓舌，擅生是非，以迷天下之主，使天下学士不反其本，妄
作孝弟，而侥幸于封侯富贵者也。子之罪大极重，疾走归！不然，我将以子肝益昼餔之膳。”


\chapter{其它附录}
其它附录的内容可以放到这里，当然如果你愿意，可
以把这部分也放到独立的文件中，然后将其 \verb|\input| 到主文件中。
%  \input{body/appendix02}
% \end{appendix}
% \end{example}
%
% \myentry{致谢声明}
% \DescribeEnv{ack}
% 把致谢做成一个环境更好一些，直接往里面写感谢的话就可以了。多希望每个人都能写上例子里的这一句啊！{\LARGE\Smiley}
% \begin{example}
% \begin{ack}
%   ……
%   还要特别感谢计算机学院胡卫谊同学在论文格式和 \LaTeX{} 编译等方面给我的很多帮助！\whutthesis让我的生活更美好！
% \end{ack}
% \end{example}
%
% \myentry{列表环境}
% \DescribeEnv{itemize}
% \DescribeEnv{enumerate}
% \DescribeEnv{description}
% 为了适合中文习惯，模板将这三个常用的列表环境用 \pkg{paralist} 对应的压缩环境替
% 换。一方面满足了多余空间的清楚，另一方面可以自己指定标签的样式和符号。细节请参
% 看 \pkg{paralist} 文档，此处不再赘述。
%
%
% \subsubsection{一些特别的命令}
% \label{ssec:special}
% \myentry{写作时标注}
% \DescribeMacro{\mynote}用来植入写作时的批注，该批注在\textsf{composing}选项下才可见，
% 提交最终版本时去掉\textsf{composing}选项，所有批注都不显示。
% \begin{example}
% 倒是作者分析到湘西的苗人的“仇外”避汉，这种思想其实并非是针对民族的，
% \mynote{为什么时阶层性的呢？这里有待查资料}{而是阶层性的}，
% 因为当时官吏的主体无疑是统治阶级的民族――这种历史背景，导致了民族性质的被强调。
% \end{example}
%
% \subsection{数学环境}
% \label{sec:math}
% \whutthesis{} 定义了常用的数学环境：
%
% \begin{center}
% \begin{tabular}{*{7}{l}}\hline
%   axiom & theorem & definition & proposition & lemma & conjecture &\\
%   公理 & 定理 & 定义 & 命题 & 引理 & 猜想 &\\\hline
%   proof & corollary & example & exercise & assumption & remark & problem \\
%   证明 & 推论 & 例子& 练习 & 假设 & 注释 & 问题\\\hline
% \end{tabular}
% \end{center}
%
% 比如：
% \begin{example}
% \begin{definition}
% 道千乘之国，敬事而信，节用而爱人，使民以时。
% \end{definition}
% \end{example}
% 产生（自动编号）：\\[5pt]
% \fbox{{\heiti 定义~1.1~~~} {道千乘之国，敬事而信，节用而爱人，使民以时。}}
%
% 列举出来的数学环境毕竟是有限的，如果想用{\heiti 胡说}这样的数学环境，那么很容易定义：
% \begin{example}
% \newtheorem{nonsense}{胡说}[chapter]
% \end{example}
%
% 然后这样使用：
% \begin{example}
% \begin{nonsense}
% 契丹武士要来中原夺武林秘笈。\pozhehao 慕容博
% \end{nonsense}
% \end{example}
% 产生（自动编号）：\\[5pt]
% \fbox{{\heiti 胡说~1.1~~~} {契丹武士要来中原夺武林秘笈。\kern0.3ex\rule[0.8ex]{2em}{0.1ex}\kern0.3ex 慕容博}}
%
% \subsection{自定义以及其它}
% \label{sec:othercmd}
% 模板的配置文件 WHUTthesis.cfg 中定义了很多固定词汇，一般无须修改。如果有特殊需求，
% 推荐在导言区使用 \cs{renewcommand}。当然，导言区里可以直接使用中文。
%
%
%\section{致谢}
%\label{sec:thanks}
%
% 在模板的开发过程中，我的女朋友给了我最大的支持，正是由于她的鼓励让我在遭遇似乎看不到尽头的困难时
% 能够坚持下来。冉冉，谢谢你！\Heart
%
% \whutthesis 的开发参考了前人的大量工作，这里特别对ThuThesis的薛瑞尼、nudtpaper
% 的刘本源、SEUThesis的许元、SJTU Master Thesis的Wei Jianwen致以最诚挚的谢意！是
% 你们无私的工作帮助我完成了这份模板，谢谢！
%
% 感谢所有在论文致谢中提及 \whutthesis{} 或者本人的同学，你们的肯定是我最大的收获！
% 在我离开后，这份模板计划由计算机
% 研学生会学术部继续维护，这里我呼唤感兴趣的同学能出手相助，给模板的开发和维护注入新的
% 活力，让我们一起把 \whutthesis{} 做得更好！
%
%
%\section{历史}
%\label{sec:history}
%
% \subsection{官方认证及推广}
% \label{ssec:auth}
% 由于论文需要经过反剽窃系统检查和最终在图书馆提交电子版两道程序，且这些系统基本只支持doc格式，
% 由此给该模板的推广带来巨大阻力。2010年03月本模板初步完成时本人去马区找过研究生院
% 主管学位论文事宜的老师以期让该模板获得官方认证，由于主管老师工作繁忙，
% 本人在余区和马区间辗转两次后此事便再无下文。提交反剽窃检查时先提交的pdf版，期待该检查
% 系统如其所声称的那样支持pdf，不幸的是奇迹没有出现，论文被打回，用Acrobat Pro 9导出一个无排版的doc
% 版应付过去。图书馆的电子版提交与上述遭遇相同，同样用Acrobat Pro 9导出的无排版的doc应付过去。真是
% 阴魂不散的doc啊……
%
% 由于未能获得官方认证，本模板并未做大力推广，目前只在googlecode建立了托管项目并在经纬论坛
% 发表了一篇发布帖(\url{http://forum.wutnews.net/bbs/topic.aspx?tid=370363})，该贴的置顶申请
% 被版主忽略。从现状来看，官方渠道仅收到学院研工办老师对该项目表示的支持。希望后续维护该项目的
% 同学能够在官方认证上取得突破并在推广上有所作为，让更多的理工学子在撰写毕业论文时受益于
% \whutthesis{}，{\heiti 拜托了！}
%
% \subsection{版本历史及维护团队}
% \label{ssec:verhis}
% \subsubsection{2010年03月-\ 2010年06月}
%       \noindent {\bf 维\qquad 护：}胡卫谊\\
%       {\bf 特别感谢：}所有测试和使用\whutthesis 的同学\\
%       {\bf 发布序列：}
%       \begin{itemize}
%           \item v0.1.2d 2010/06/19 文档更新
%               \begin{enumerate}
%                   \item 补充todolist，增加“\ref{ssec:auth}官方认证及推广”小节
%               \end{enumerate}
%           \item v0.1.2c 2010/04/19 新特性和功能修复
%               \begin{enumerate}
%                   \item 引入footmisc宏包以修改脚注样式 
%                   \item 引入todonotes宏包以支持编写时的批注。非常漂亮的批注，需要新增的composing文档选项的支持，具体使用方法请看用户手册 
%                   \item 编写论文时输出文稿背景色设置为橙色新闻纸色，保护视力 
%                   \item 子图引用加上小括号，如图 1-2(a)所示 
%                   \item 将参考文献样式文件转换为utf-8编码，正确地生成文献列表 
%                   \item 编译脚本的 msmake my 选项引入第二个选项once，只处理文档一次以加快编译速度，方便即时查看 
%               \end{enumerate}
%           \item v0.1.2b 2010/03/31 微小调整。私房版，未发布。
%               \begin{enumerate}
%                   \item 引入footmisc宏包以修改脚注样式
%                   \item 引入todonotes宏包以支持编写时的批注
%                   \item 用户手册增加版本历史和TODOlist章节
%               \end{enumerate}
%           \item v0.1.2a 2010/03/29 微小调整
%               \begin{enumerate}
%                   \item 发布版本号后加字母后缀，避免微小调整造成版本号上升过快 
%                   \item 调整了部分宏包的加载顺序，脚注的超链接得以修复 
%                   \item 取消了扉页英文题目的悬挂缩进，英文题目的第二行下划线与第一行对齐 
%               \end{enumerate}
%           \item v0.1.1 2010/03/22 微小调整
%               \begin{enumerate}
%                   \item 完善了数学字体，使用了与同济大学第五版高数教材所用字体相同的字体，使得数学公式更加漂亮 
%                   \item 压缩了列表项间距，中文版式下更加漂亮 
%               \end{enumerate}
%           \item v0.1.0 2010/03/15 初始发布
%               \begin{enumerate}
%                   \item 该版本目前只完整支持硕士论文，博士论文的双面未调试
%               \end{enumerate}
%       \end{itemize}
% \subsubsection{2010年07月-\ ……}
%       \noindent {\bf 维\qquad 护：}\\
%       {\bf 特别感谢：}\\
%       {\bf 发布序列：}
%
% \subsection{TODOlist}
% \label{ssec:todo}
% 下表记录了计划实现的功能和特性
%   \begin{enumerate}
%       \item 调试博士论文的双面支持
%       \item 加入本科论文的支持
%       \item 修正“致谢”部分页眉风格的dirty code
%       \item 增加编写时的批注功能\qquad  (DONE! v0.1.2b)
%       \item 去掉图表标题后的冒号（这个问题尝试了许多办法均未能解决）
%       \item 修改listings宏包的标题序号中的分隔符为'-'
%       \item 调整章开始页上部空白大小
%       \item 调整a3封面的排版
%       \item 修正PDF文件中从目录点击超链接到“参考文献”时目的页不匹配问题
%       \item 将“目录”插入pdf的书签中，方便导航
%       \item 去掉|\mynote{}|与前文间的一段小间距
%       \item Linux平台迁移及测试
%   \end{enumerate}
%
% \StopEventually{\PrintChanges\PrintIndex}
% \clearpage
%
% \section{实现细节}
%
% \subsection{基本信息}
%    \begin{macrocode}
%<cls>\NeedsTeXFormat{LaTeX2e}[1999/12/01]
%<cls>\ProvidesClass{WHUTthesis}
%<cfg>\ProvidesFile{WHUTthesis.cfg}
%<cls|cfg>[2010/03/02 v0.1.0 Thesis Template for Wuhan University of Technology]
%    \end{macrocode}
%
% \subsection{定义选项}
% \label{sec:defoption}
%<*cls>%
%    \begin{macrocode}
\def\whutthesis{\textsc{WHUT}\textsc{thesis}}
%%\def\version{0.1.2b}
\newif\ifwhut@master\whut@masterfalse
\newif\ifwhut@doctor\whut@doctorfalse
\newif\ifwhut@istwoside\whut@istwosidefalse
\newif\ifwhut@composing\whut@composingfalse
\DeclareOption{twoside}{\whut@istwosidetrue}
\DeclareOption{master}{\whut@mastertrue}
\DeclareOption{doctor}{\whut@doctortrue}
\DeclareOption{composing}{\whut@composingtrue}
\DeclareOption*{\PackageWarning{whutthesis}{Unknown Option '\CurrentOption'}}
%    \end{macrocode}
%
%% 将选项传递给 ctexbook 类
%% 用户至少要提供一个选项：指定论文类型。
%    \begin{macrocode}
\DeclareOption*{\PassOptionsToClass{\CurrentOption}{ctexbook}}
\ProcessOptions\relax
\ifwhut@istwoside
    \LoadClass[cs4size,a4paper,twoside,openright,fancyhdr,fntef]{ctexbook}
\else
    \LoadClass[cs4size,a4paper,oneside,openany,fancyhdr,fntef]{ctexbook}
\fi
\ifwhut@master\relax\else
    \ifwhut@doctor\relax\else
    \ClassError{whutthesis}%
        {You have to specify one of thesis options: master or doctor.}{}
    \fi
\fi
%    \end{macrocode}
%
% \subsection{装载宏包}
% \label{sec:loadpackage}
%% 首先调用在文档类书写中需要的过程控制语句
%    \begin{macrocode}
\RequirePackage{ifthen,calc}
%    \end{macrocode}
%
%% 我们直接用\textsf{geometry}宏包进行页面边距的设定，调用titlesec设定标题以及页眉页脚，
%% 用\textsf{titletoc}设定目录格式。
%    \begin{macrocode}
\RequirePackage{geometry}
%%\RequirePackage[center]{titlesec} %此包会覆盖ctexbook的设置，取消掉
\RequirePackage{titletoc}
%    \end{macrocode}
%
%% 图形支持宏包。
% 文档中另外重要的两个部分是表格和图片，首先来看图片：\textsf{graphicx}宏包是必不可少的，
% 并排图形。\textsf{subfigure} 已经不再推荐，用新的 \textsf{subfig}。加入 \verb|config| 选项
% 以便兼容 \textsf{subfigure} 的命令。浮动图形和表格标题样式。\textsf{caption2} 已经不
% 推荐使用，采用新的 \textsf{caption}。它会自动被 \textsf{subfig} 装载进来。所以可以在
% 后面看到 \textbf{captionsetup} 的命令，后一个包\textsf{float}的作用是可以用H命令，
% 将浮动对象强制放在这里(副作用是版面可能不好)：
%    \begin{macrocode}
\RequirePackage{graphicx}
\RequirePackage[config]{subfig}
\RequirePackage{float}
%    \end{macrocode}
%
% 表格中，我们采用\textsf{longtable}来处理长的表格，还需要\textsf{array}包，
% 标准的论文需要表格为三线表，这里引用\textsf{booktabs}宏包来处理，
% 这样，我们就可以使用\verb|\toprule|，\verb|\midrule|，\verb|bottomrulle|这样的命令，
% 在表格中支持跨行，需要引入\textsf{multirow}包，\textsf{tabularx}的作用是为了引入
% 固定宽度的表格，\textsf{slashbox}可以让我们在表格中使用反斜线：
%    \begin{macrocode}
\RequirePackage{array}
\RequirePackage{longtable}
\RequirePackage{booktabs}
\RequirePackage{multirow}
%%\RequirePackage{tabularx} %\textsf{tabularx}需要在\textsf{hyperref}之后引入才不会使正文的footnote的超链接失效
\RequirePackage{slashbox}
%    \end{macrocode}
%
%% 引入\textsf{paralist}来达到比较好看的列表环境
% \changes{v0.1.1}{2010/03/22}{加入列表行距的设置，来自thuthesis}
%    \begin{macrocode}
\RequirePackage[neverdecrease]{paralist}
\setdefaultenum{1)}{a)}{i.}{A.}% 设置列表模式
\setdefaultleftmargin{4ex}{}{}{}{}{}
\let\itemize\compactitem
\let\enditemize\endcompactitem
\let\enumerate\compactenum
\let\endenumerate\endcompactenum
\let\description\compactdesc
\let\enddescription\endcompactdesc
%    \end{macrocode}
% 
%% 文档中还需要一定的色彩控制和字体控制
% \changes{v0.1.2c}{2010/04/15}{论文编写模式下的文档背景色，新闻纸的淡黄色，保护视力}
%    \begin{macrocode}
\RequirePackage{xcolor}
\ifwhut@composing
    \pagecolor[rgb]{1,0.95,0.88}
\fi
%    \end{macrocode}

% 为了排出漂亮的数学公式，\textsf{amsmath}包是必不可少的，\textsf{txfonts}的作用是用
% 自己的typewriter字体替换系统Courier字体，它必须在\AmSTeX{}之后，数学应用中还需要
% 定理环境，我们一并包括进来：
% \changes{v0.1.1}{2010/03/22}{去掉txfonts包的varg选项，使得字体更接近同济大学第五版高数和thuthesis中的字体设置}
% \changes{v0.1.1}{2010/03/22}{加载txfonts后会覆盖掉ctexbook加载的fontspec包的英文字体设置，加入重设这些字体的修正}
%    \begin{macrocode}
%%\RequirePackage{amsmath,amsfonts,amssymb,bm,mathrsfs}
\RequirePackage{amsmath,amssymb}
\RequirePackage{txfonts}
\setmainfont{Times New Roman}
\setsansfont{Arial}
\setmonofont{Courier New}
\RequirePackage[amsmath,thmmarks,hyperref]{ntheorem}
%    \end{macrocode}
%
%% 图表双标题支持。和其他包（可能是subfig）冲突！取消掉{\LARGE\Frowny}
%    \begin{macrocode}
%%\RequirePackage{ccaption}
%%\captiondelim{\ }
%%\captionnamefont{\kaishu\zihao{5}}
%%\captiontitlefont{\kaishu\zihao{5}}
%    \end{macrocode}
%
%% 将参考文献、索引、表格目录和图形
%% 目录等标题及其所在页码自动地添加到标
%% 题目录，可以断行（用\protect\\）@@TODO:来自seu，有待测试
%    \begin{macrocode}
%%\RequirePackage{makeidx} % 索引  
%%\RequirePackage[nottoc]{tocbibind}
%    \end{macrocode}
%
%% 定公式、图、表编号为"3-1"的形式，即分隔符由.变为短杠。
%% 子图和子表的应用序号外加小括号，如图 1-2(a)。
% \changes{v0.1.2c}{2010/04/15}{子表引用的时候子序号没有外面的小括号1-2(a)，加上}
%    \begin{macrocode}
\renewcommand\theequation{\arabic{chapter}--\arabic{equation}}
\renewcommand\thefigure{\arabic{chapter}--\arabic{figure}}
\renewcommand\thetable{\arabic{chapter}--\arabic{table}}
\captionsetup[table]{position=top,belowskip={12bp-\intextsep},aboveskip=6bp}
\captionsetup[figure]{position=bottom,belowskip={12bp-\intextsep},aboveskip=6bp}
\captionsetup[subfloat]{labelformat=simple,captionskip=6bp,nearskip=6bp,farskip=0bp,topadjust=0bp}
\renewcommand{\thesubfigure}{(\alph{subfigure})}
\renewcommand{\thesubtable}{(\alph{subtable})}
%    \end{macrocode}
%
%% 首行缩进宏包
%    \begin{macrocode}
\RequirePackage{indentfirst}
%    \end{macrocode}
%
%% 如果插入的图片没有指定扩展名，那么依次搜索下面的扩展名所对应的文件
%    \begin{macrocode}
\DeclareGraphicsExtensions{.pdf,.eps,.jpg,.png}
%    \end{macrocode}
%
%
% \changes{v0.1.2a}{2010/03/29}{\textsf{setspace}宏包提到hyperref前加载，避免脚注超链接失效}
%% 行距,学校要求固定值20pt，每页可打印正文约32行
%% |\linestrech| 的值为 1 时为单倍行距， 1.2 时是一倍半行距， 而为 1.6 时是双倍行距。
%%其实不同尺寸的字体行间距都不相同，而是成比例关系。这个20pt是对正文主要字体来说的。
%%在TeX中基本的行间距是|\baselineskip|, 对于12pt的字体，这个值等于14.5pt, 而真正的行间距是
%%|\baselineskip * \baselinestretch| ，|\baselinestretch|默认为1, 但我们可以重新设置它的值，如 
%%|\renewcommand{\baselinestretch}{1.38}| 就得到真正的行间距为14.5pt*1.38≈20pt。
%%而这样定义之后，对不同尺寸的字体都会按同样的比例因子1.38放大行间距，使得全文排版能协调一致。
%    \begin{macrocode}
\RequirePackage{setspace}
\setstretch{1.38}
%    \end{macrocode}
%
% \changes{v0.1.2b}{2010/03/30}{引入\textsf{footmisc}宏包以修改脚注样式}
%% \textsf{footmisc}宏包可以方便地更改脚注样式。这里修改了宏包默认的bringhurs样式，
%% \textbf{注意：此脚注样式只支持每页18个脚注，多于18个将无法编号}
%    \begin{macrocode}
\RequirePackage[perpage,symbol*,stable]{footmisc}
\DefineFNsymbolsTM*{hurst}{%
	\textdagger       \dagger
      \textdaggerdbl  \ddagger
      \textsection      \mathsection
      \textparagraph  \mathparagraph
      \textbardbl       \|%
      {\textasteriskcentered\textasteriskcentered}{**}%
      {\textdagger\textdagger}{\dagger\dagger}%
      {\textdaggerdbl\textdaggerdbl}{\dagger\dagger}%
      {\textsection\textsection}{\mathsection\mathsection}%
      {\textparagraph\textparagraph}{\mathparagraph\mathparagraph}%
      {\textasteriskcentered\textasteriskcentered\textasteriskcentered}{***}%
      {\textdagger\textdagger\textdagger}{\dagger\dagger\dagger}%
      {\textdaggerdbl\textdaggerdbl\textdaggerdbl}{\ddagger\ddagger\ddagger}%
      {\textsection\textsection\textsection}%%
      {\mathsection\mathsection\mathsection}%
      {\textparagraph\textparagraph\textparagraph}%%
      {\mathparagraph\mathparagraph\mathparagraph}%
      \textasteriskcentered  *
}%
\setfnsymbol{hurst}
%    \end{macrocode}
%
% \begin{macro}{\mynote}
% \changes{v0.1.2b}{2010/03/30}{引入\textsf{todonotes}宏包以支持编写论文时自己写批注，批注的可见性由composing文档选项控制}
%% 引入\textsf{todonotes}宏包以支持编写论文时自己写批注，批注的可见性由composing文档选项控制。
%% 也就是定义了composing文档选项时批注可见，不定义时不可见。
%% 自定义\textsf{mynote}命令，以方便对选定文字做批注。
%% 引入\textsf{etex}宏包以避免\textsf{todonotes}和\textsf{listings}宏包冲突
%% 而导致 \textsf{No room for a new count}错误
%    \begin{macrocode}
\newcommand\mynote[2]{
    {#2}}
\ifwhut@composing
    \RequirePackage{etex}
    \RequirePackage[textwidth=2.9cm,shadow, backgroundcolor=lime]{todonotes}
    \renewcommand\mynote[2]{
        \todo{#1}
        {\color{gray}#2}}
\fi
%    \end{macrocode}
% \end{macro}


% 另外一个关键部分是文献索引，包括书签以及参考文献的索引，记得\textsf{hyperref}配合
% \XeTeX{}使用时不能开启Unicode选项，而\textsf{hypernat}可以让\textsf{hyperref}和
% \textsf{natbib}混合使用，需要放在这两个之后：
%    \begin{macrocode}
\RequirePackage[numbers,super,sort&compress]{natbib}
\RequirePackage{hyperref}
\hypersetup{%
    CJKbookmarks=true,
    bookmarksnumbered=true,
    bookmarksopen=true,
    bookmarksopenlevel=0,
    breaklinks=true,
    colorlinks=false,
    plainpages=false,
    pdfpagelabels,
    pdfborder=0 0 0}
\RequirePackage{hypernat}
%    \end{macrocode}

% \changes{v0.1.2a}{2010/03/29}{延后到hyperref加载后再加载tabularx，避免脚注的超链接失效}
%% \textsf{tabularx}需要在\textsf{hyperref}之后引入才不会使正文的footnote的超链接失效
%    \begin{macrocode}
\RequirePackage{tabularx}
%    \end{macrocode}

%% 设置 url 样式，与上下文一致
%    \begin{macrocode}
\urlstyle{same}
%    \end{macrocode}
%
%% 下划线自动换行
%    \begin{macrocode}
\RequirePackage{ulem}
%    \end{macrocode}
%
%% 设置章节格式
%    \begin{macrocode}
\CTEXsetup[number={\arabic{chapter}},name={第,章},
            nameformat={\bfseries\centering\zihao{-2}},
            titleformat={\bfseries\zihao{-2}},
            afterskip={30pt}]{chapter}
\CTEXsetup[nameformat={\bfseries\zihao{3}},
            titleformat={\bfseries\zihao{3}}]{section}
\CTEXsetup[nameformat={\bfseries\zihao{4}},
            titleformat={\bfseries\zihao{4}}]{subsection}
\CTEXsetup[nameformat={\bfseries\zihao{-4}},
            titleformat={\bfseries\zihao{-4}}]{subsubsection}            
\CTEXsetup[format={\Large\bfseries}]{section}
\CTEXsetup[beforeskip={10pt}]{chapter}
%    \end{macrocode}
%</cls>
%
%
% \subsection{字符串常量定义}
% \label{sec:strdef}
%    \begin{macrocode}
%<*cfg>
% \theoremsymbol{\ensuremath{\blacksquare}}
\theoremsymbol{\ensuremath{\square}}
%\theoremstyle{nonumberplain}
\newtheorem*{proof}{证明}
\theoremstyle{plain}
\theoremsymbol{}
\theoremseparator{：}
\newtheorem{assumption}{假设}[chapter]
\newtheorem{definition}{定义}[chapter]
\newtheorem{proposition}{命题}[chapter]
\newtheorem{lemma}{引理}[chapter]
\newtheorem{theorem}{定理}[chapter]
\newtheorem{axiom}{公理}[chapter]
\newtheorem{corollary}{推论}[chapter]
\newtheorem{exercise}{练习}[chapter]
\newtheorem{example}{例}[chapter]
\newtheorem{remark}{注释}[chapter]
\newtheorem{problem}{问题}[chapter]
\newtheorem{conjecture}{猜想}[chapter]
\renewcommand\contentsname{目\hspace{1em}录}
\renewcommand\listfigurename{插图索引}
\renewcommand\listtablename{表格索引}
\newcommand\listequationname{公式索引}
\newcommand\equationname{公式}
\renewcommand\bibname{参考文献}
\renewcommand\indexname{索引}
\renewcommand\figurename{图}
\renewcommand\tablename{表}
\newcommand{\cabstractname}{摘\hspace{1em}要}
\newcommand{\eabstractname}{Abstract}
\let\CJK@todaysave=\today
\def\CJK@todaysmall@short{\the\year 年 \the\month 月}
\def\CJK@todaysmall{\CJK@todaysmall@short \the\day 日}
\CTEXdigits{\whut@yeartmp}{\the\year}
\CTEXnumber{\whut@monthtmp}{\the\month}
\def\CJK@todaybig@short{\whut@yeartmp \ 年 \whut@monthtmp \ 月}
\def\CJK@todaybig{\CJK@todaybig@short \CJKnumber{\the\day} 日}
\def\CJK@today{\CJK@todaysmall}
\renewcommand\today{\CJK@today}
\newcommand\CJKtoday[1][1]{%
  \ifcase#1\def\CJK@today{\CJK@todaysave}
    \or\def\CJK@today{\CJK@todaysmall}
    \or\def\CJK@today{\CJK@todaybig}
  \fi}
\cdate{\CJK@todaybig@short}
\edate{\ifcase \month \or January\or February\or March\or April\or May%
		\or June\or July \or August\or September\or October\or November
		\or December\fi\unskip,\ \ \the\year}
%% 封面等一些标题串
\def\whut@title@sep{：}
\def\whut@ckeywords@separator{，}
\def\whut@ekeywords@separator{,}
\def\whut@categorynumbercap{分类号}
\def\whut@secretlevelcap{密级} %% 注明论文密级为公开、内部、秘密或机密
\def\whut@UDCcap{UDC}
\def\whut@schoolcodecap{学校代码}
\def\whut@schoolcode{10497}
\def\whut@thesisname{学位论文}
\def\whut@ctitlecap{题目}
\def\whut@etitlecap{英文题目}
\def\whut@authornamecap{研究生姓名}
\def\whut@advisorcap{指导教师}
\def\whut@advisornamecap{姓名}
\def\whut@advisortitlecap{职称}
\def\whut@advisordegreecap{学位}
\def\whut@advisordepartmentcap{单位名称}
\def\whut@advisorzipcodecap{邮编}
\def\whut@coadvisorcap{副指导教师}
\def\whut@applycap{申请学位级别}
\def\whut@majorcap{学科专业名称}
\def\whut@submitdatecap{论文提交日期}
\def\whut@oraldefensedatecap{论文答辩日期}
\def\whut@authorizeorganizationcap{学位授予单位}
\def\whut@authorizedatecap{学位授予日期}
\def\whut@chairmancap{答辩委员会主席}
\def\whut@annotatorcap{评阅人}
\def\whut@title@sep{：}
\def\whut@secretlevel{秘密}
\def\whut@secretyear{\the\year}
\def\whut@schoolname{武汉理工大学}
\def\whut@apply{（申请\whut@cdegree 学位论文）}
\def\whut@cdepartmentcap{培养单位}
\def\whut@cmajortitle{学科专业}
\def\whut@cauthortitle{研究生}
\def\whut@cadvisortitle{指导教师}
\def\whut@masterheader{武汉理工大学硕士学位论文}
\def\whut@doctorheader{武汉理工大学博士学位论文}
\newcommand{\whut@declaretitle}{独创性声明}
\newcommand{\whut@declare}{%
本人声明，所呈交的论文是本人在导师指导下进行的研究工作及取得的研究成果。
尽我所知，除了文中特别加以标注和致谢的地方外，论文中不包含其他人已经发表
或撰写过的研究成果，也不包含为获得武汉理工大学或其他教育机构的学位或证书而使用过的材料。
与我一同工作的同志对本研究所做的任何贡献均已在论文中作了明确的说明并表示了谢意。
}
\def\whut@declaresigncap{签名：}
\def\whut@declaredatecap{日期：}
\newcommand{\whut@authorizationtitle}{学位论文使用授权书}
\newcommand{\whut@authorization}{%
本人完全了解武汉理工大学有关保留、使用学位论文的规定，
即学校有权保留并向国家有关部门或机构送交论文的复印件和电子版，
允许论文被查阅和借阅。
本人授权武汉理工大学可以将本学位论文的全部内容编入有关数据库进行检索，
可以采用影印、缩印或其他复制手段保存或汇编本学位论文。
同时授权经武汉理工大学认可的国家有关机构或论文数据库使用或收录本学位论文，
并向社会公众提供信息服务。
}
\newcommand{\whut@authorizationaddon}{（保密的论文在解密后应遵守此规定）}
\newcommand{\whut@authorsig}{研究生（签名）：}
\newcommand{\whut@teachersig}{导师（签名）：}
\newcommand{\whut@frontdate}{日期：}
\newcommand{\whut@ckeywords@title}{关键词：}
\newcommand{\whut@ekeywords@title}{Key words:}
\newcommand{\whut@denotation@name}{主要符号对照表}
\newcommand{\whut@ackshortname}{致谢}
\newcommand{\whut@ackname}{致\hspace{1em}谢}
\newcommand{\whut@resume@title}{个人简历、在学期间发表的学术论文与研究成果}
%</cfg>
%    \end{macrocode}
%
%
% \subsection{主文档格式}
% \label{sec:mainbody}
%
% \subsubsection{页幅设置}
% \label{sec:layout}
%<*cls>
%% 正文统一用小四号字，行间距：固定值20磅，页眉：2.6cm，页脚：2.4cm，
%% 论文用A4纸标准大小的白纸打印，页边距上下：3.5cm，左右：3.2cm，四周适当修边即可。
%    \begin{macrocode}
\geometry{top=3.5cm,bottom=3.5cm,left=3.2cm,right=3.2cm}
\geometry{headheight=2.6cm,headsep=3mm,footskip=13mm}
%    \end{macrocode}
%</cls>
%
% \subsubsection{页眉页脚}
% \label{sec:headerfooter}
%    \begin{macrocode}
%<*cls>
%% 页眉页脚
\pagestyle{fancy}
\fancyhf{}
\ifwhut@master
    \fancyhead[CE,CO]{\small \whut@masterheader}
\else
    \fancyhead[CE,CO]{\small \whut@doctorheader}
\fi
% \fancyfoot[OL]{\bfseries \small ~---~\thepage~---~}
\fancyfoot[C]{\small ~---~{\bf\thepage}~---~}      % 页脚
% \fancyfoot[LE,RO]{\bfseries \small ~--~\thepage~--~}
\fancypagestyle{plain}{% 设置开章页页眉页脚风格
    \fancyhf{}%
    % \fancyfoot[C]{\bfseries \small ~第~\thepage~页~}
    \fancyfoot[C]{\small ~---~{\bf\thepage}~---~} %%首页页脚格式
    \renewcommand{\headrulewidth}{0pt}%
    \renewcommand{\footrulewidth}{0pt}%
}
%</cls>
%    \end{macrocode}
%
%<*cls>
%\subsubsection{前置目录}
%% 前置部分的封面在后面详细介绍，首先看目录，要求为：目次页由论文的章、节、条、项、附录等的序号、名称和页码组成，
%% 另页排在序之后。目次页标注学位论文的前三级目录。标题统一用“目录”，黑体3字号字居中，段前、段后间距为1行;
%% 各章（一级目录）名称用黑体小4号字，段前间距为0.5行，段后间距为0行; 其它（二、三级目录）用宋体小4号字，
%% 段前、段后间距为0行。：
%% 在\LaTeX{}中，Chapter在目录中默认是没有点的,我们加上,另外我们一并将目录中的section和subsection设定好,
%
%    \begin{macrocode}
\titlecontents{chapter}[0pt]{\vspace{0.25\baselineskip} \heiti \zihao{-4}}
    {\thecontentslabel\quad}{}
    {\hspace{.5em}\titlerule*{.}\contentspage}
\titlecontents{section}[2em]{\songti \zihao{-4}}
    {\thecontentslabel\quad}{}
    {\hspace{.5em}\titlerule*{.}\contentspage}
\titlecontents{subsection}[4em]{\songti \zihao{-4}}
    {\thecontentslabel\quad}{}
    {\hspace{.5em}\titlerule*{.}\contentspage}
%    \end{macrocode}
%</cls>
%
%
% \subsubsection{封面和封底}
% \label{sec:cover}
% \begin{macro}{\whut@define@term}
%    \begin{macrocode}
%<*cls>
%% 方便的定义封面的一些替换命令。
\def\whut@define@term#1{
    \expandafter\gdef\csname #1\endcsname##1{%
    \expandafter\gdef\csname whut@#1\endcsname{##1}}
    \csname #1\endcsname{}}
%</cls>
%    \end{macrocode}
% \end{macro}
%
% \begin{macro}{\cauthordepartment}
% \begin{macro}{\categorynumber}
% \begin{macro}{\secretlevel}
% \begin{macro}{\UDC}
% \begin{macro}{\ctitle}
% \begin{macro}{\ctitleformated}
% \begin{macro}{\etitle}
% \begin{macro}{\etitleformated}
% \begin{macro}{\cauthor}
% \begin{macro}{\cadvisorname}
% \begin{macro}{\cadvisortitle}
% \begin{macro}{\cadvisordegree}
% \begin{macro}{\cadvisordepartment}
% \begin{macro}{\advisorzipcode}
% \begin{macro}{\ccoadvisorname}
% \begin{macro}{\ccoadvisortitle}
% \begin{macro}{\ccoadvisordegree}
% \begin{macro}{\ccoadvisordepartment}
% \begin{macro}{\coadvisorzipcode}
% \begin{macro}{\cmajor}
% \begin{macro}{\csubmitdate}
% \begin{macro}{\coraldefensedate}
% \begin{macro}{\cauthorizedepartment}
% \begin{macro}{\cauthorizedate}
% \begin{macro}{\cchairman}
% \begin{macro}{\cannotatora}
% \begin{macro}{\cannotatorb}
% \begin{macro}{\cdate}
% \begin{macro}{\cdegree}
% \begin{macro}{\csubject}
% \begin{macro}{\edegree}
% \begin{macro}{\edepartment}
% \begin{macro}{\emajor}
% \begin{macro}{\esubject}
% \begin{macro}{\eauthor}
% \begin{macro}{\eadvisorname}
% \begin{macro}{\eassosupervisor}
% \begin{macro}{\ecoadvisorname}
% \begin{macro}{\edate}
%% 封面、摘要、版权、致谢变量定义
%    \begin{macrocode}
%<*cls>
\newcommand\cauthordepartment[2][]{\def\whut@cauthordepartment@short{#1}\def\whut@cauthordepartment{#2}}
\whut@define@term{categorynumber}
\whut@define@term{secretlevel}
\whut@define@term{UDC}
\whut@define@term{ctitle}
\whut@define@term{ctitleformated}
\whut@define@term{etitle}
\whut@define@term{etitleformated}
\whut@define@term{cauthor}
\whut@define@term{cadvisorname}
\whut@define@term{cadvisortitle}
\whut@define@term{cadvisordegree}
\whut@define@term{cadvisordepartment}
\whut@define@term{advisorzipcode}
\whut@define@term{ccoadvisorname}
\whut@define@term{ccoadvisortitle}
\whut@define@term{ccoadvisordegree}
\whut@define@term{ccoadvisordepartment}
\whut@define@term{coadvisorzipcode}
\whut@define@term{cmajor}
\whut@define@term{csubmitdate}
\whut@define@term{coraldefensedate}
\whut@define@term{cauthorizedepartment}
\whut@define@term{cauthorizedate}
\whut@define@term{cchairman}
\whut@define@term{cannotatora}
\whut@define@term{cannotatorb}
\whut@define@term{cdate}
\whut@define@term{cdegree}
\def\whut@cauthordepartment@short{}
\def\whut@cauthordepartment{}
\whut@define@term{edegree}
\whut@define@term{edepartment}
\whut@define@term{emajor}
\whut@define@term{eauthor}
\whut@define@term{eadvisorname}
\whut@define@term{ecoadvisorname}
\whut@define@term{edate}
%</cls>
%    \end{macrocode}
% \end{macro}
% \end{macro}
% \end{macro}
% \end{macro}
% \end{macro}
% \end{macro}
% \end{macro}
% \end{macro}
% \end{macro}
% \end{macro}
% \end{macro}
% \end{macro}
% \end{macro}
% \end{macro}
% \end{macro}
% \end{macro}
% \end{macro}
% \end{macro}
% \end{macro}
% \end{macro}
% \end{macro}
% \end{macro}
% \end{macro}
% \end{macro}
% \end{macro}
% \end{macro}
% \end{macro}
% \end{macro}
% \end{macro}
% \end{macro}
% \end{macro}
% \end{macro}
% \end{macro}
% \end{macro}
% \end{macro}
% \end{macro}
% \end{macro}
% \end{macro}
% \end{macro}
%
%
%
% \begin{macro}{\whut@first@titlepage}
% 封面和封二！
% \changes{v0.1.2a}{2010/03/29}{取消了英文题目第二行的缩进，保持与第一行下划线对齐}
%    \begin{macrocode}
%<*cls>
%%
%%
%% 封面
\newlength{\whut@title@width}
\def\whut@put@covertitle#1{\makebox{\hb@xt@ \whut@title@width{#1}}}
%%封面
\newcommand{\whut@first@coverpage}{
    \begin{center}
        \vspace*{-1.5cm}
%%      绘制学校名
        \begin{picture}(0,0)(115,110)
            {\includegraphics[width=0.56\textwidth]{figure/whut.pdf}}
        \end{picture}\\
%%      绘制论文题目
        \vspace*{4.3cm}
        \zihao{4}{\whut@apply} \\
        \vspace*{3.2cm}
        \heiti\zihao{1}\whut@ctitle \\
%%      绘制作者等信息
        \vspace*{3cm}
        \setlength{\whut@title@width}{3em}
        {\bf\zihao{3}\fangsong
        \begin{tabular}{p{\whut@title@width}@{\extracolsep{4pt}}l@{\extracolsep{2pt}}l}
            \whut@put@covertitle{\whut@cdepartmentcap} & \whut@title@sep & \whut@cauthordepartment \\
            \whut@put@covertitle{\whut@cmajortitle} & \whut@title@sep &\whut@cmajor \\
            \whut@put@covertitle{\whut@cauthortitle} & \whut@title@sep & \whut@cauthor \\
            \whut@put@covertitle{\whut@advisorcap} & \whut@title@sep & \whut@cadvisorname
        \end{tabular}}\\
%%      绘制日期
        \vspace*{2.5cm}{\zihao{3} \songti \whut@cdate}
    \end{center}
}
%%封二
\newcommand{\CoverTopOption}[3][s]{
    \makebox[2cm][#1]{\bf\songti\zihao{4} #2}{
    \underline{\makebox[3cm][c]{\bf\songti\zihao{4} #3}}}
}
\newcommand{\CoverBottomOption}[3][4.2cm]{
    \makebox[6.6cm]{
        {\songti\zihao{4}#2} %宋体四号
        \underline{\makebox[#1][c]{\kaishu\zihao{4}#3}}}
}
\def\whut@put@titlepagetitle#1{{\zihao{4}\songti \makebox{\hb@xt@ \whut@title@width{#1}}}}
\def\whut@put@titlepagecontent#1{{\zihao{4}\kaishu {#1}}}
\newcommand{\whut@first@titlepage}{
    \begin{center}
        \vspace*{-1.3cm}
        \renewcommand{\arraystretch}{1.6}
%%      绘制顶部分类号、学校名等信息
        \begin{tabular*}{\textwidth}{l@{\extracolsep\fill}l}
            \CoverTopOption{\whut@categorynumbercap}{\whut@categorynumber} &
            \CoverTopOption{\whut@secretlevelcap}{\whut@secretlevel}\\
            \CoverTopOption{\whut@UDCcap}{\whut@UDC} &
            \CoverTopOption{\whut@schoolcodecap}{\whut@schoolcode}
        \end{tabular*}\\
%%      绘制学校名
        \begin{picture}(0,0)(115,100)
            {\includegraphics[width=0.56\textwidth]{figure/whut.pdf}}
        \end{picture}\\
        \vspace*{3.5cm}
        \bf\zihao{1}{\ziju{1.5}\whut@thesisname}
    \end{center}
    
    \vspace*{0.9cm}
    \begin{center}
        \songti\zihao{4}
        \setlength{\whut@title@width}{5em}
        \setlength{\extrarowheight}{4pt}
%%      绘制题目、作者、指导教师、审阅等信息
        \hangafter 1
        \hangindent 5.23em
        \noindent\whut@put@titlepagetitle{\whut@ctitlecap}  \whut@put@titlepagecontent{\uline{\hfill \whut@ctitleformated \hfill}}\\
        \noindent\whut@put@titlepagetitle{\whut@etitlecap}  \uline{\hfill \whut@etitleformated \hfill}\\
        \noindent\whut@put@titlepagetitle{\whut@authornamecap} \whut@put@titlepagecontent{\uline{\hfill \whut@cauthor \hfill}} \\ \par
        \begin{tabular*}{\textwidth}{@{}p{\whut@title@width}l@{}c@{}c}
                                                                            & \whut@advisornamecap \underline{\makebox[3.7cm]{\kaishu\whut@cadvisorname}} \whut@advisortitlecap \underline{\makebox[2.4cm]{\kaishu\whut@cadvisortitle}}  & \whut@advisordegreecap \underline{\makebox[2.4cm]{\kaishu\whut@cadvisordegree}}\\
            \raisebox{1.8ex}[0pt]{\whut@put@titlepagetitle{\whut@advisorcap}} & \whut@advisordepartmentcap \underline{\makebox[6.2cm]{\kaishu\whut@cadvisordepartment}} 													& \whut@advisorzipcodecap \underline{\makebox[2.4cm]{\kaishu\whut@advisorzipcode}}
        \end{tabular*}\\
        \begin{tabular*}{\textwidth}{@{}p{\whut@title@width}l@{}c@{}c}
                                                                            & \whut@advisornamecap \underline{\makebox[3.7cm]{\kaishu\whut@ccoadvisorname}} \whut@advisortitlecap \underline{\makebox[2.4cm]{\kaishu\whut@ccoadvisortitle}}  & \whut@advisordegreecap \underline{\makebox[2.4cm]{\kaishu\whut@ccoadvisordegree}}\\
            \raisebox{1.8ex}[0pt]{\whut@put@titlepagetitle{\whut@coadvisorcap}} & \whut@advisordepartmentcap \underline{\makebox[6.2cm]{\kaishu\whut@ccoadvisordepartment}} 													& \whut@advisorzipcodecap \underline{\makebox[2.4cm]{\kaishu\whut@coadvisorzipcode}}
        \end{tabular*}\\
        \begin{tabular*}{\textwidth}{@{}ll@{}}
            \CoverBottomOption{\whut@applycap}{\whut@cdegree} & \CoverBottomOption{\whut@majorcap}{\whut@cmajor}\\
            \CoverBottomOption{\whut@submitdatecap}{\whut@csubmitdate} & \CoverBottomOption{\whut@oraldefensedatecap}{\whut@coraldefensedate}\\
            \CoverBottomOption{\whut@authorizeorganizationcap}{\whut@cauthorizedepartment} & \CoverBottomOption{\whut@authorizedatecap}{\whut@cauthorizedate}\\
%%            \CoverBottomOption[4cm]{\whut@chairmancap}{\whut@cchairman} & \CoverBottomOption{\whut@annotatorcap}{\whut@cannotatora}\\ %%这种方式始终不能对齐，不是表格左边有半个字间距就是主席行凸出去。只能把command释放出来写，dirty code!
            \CoverBottomOption[3.8cm]{\whut@chairmancap}{\whut@cchairman} & \makebox[7.2cm][s]{\songti\zihao{4}{\whut@annotatorcap}\underline{\makebox[4.3cm][c]{\kaishu\zihao{4}\whut@cannotatora}}}\\
                    & \CoverBottomOption{\makebox[2.65cm]{}}{\whut@cannotatorb}
        \end{tabular*}\\
%%      绘制论文打印日期
%%        \vspace*{3.2cm}{\zihao{3} \songti \whut@cdate}
        \vskip2.2cm{\zihao{3} \songti \whut@cdate}
    \end{center}
} % end of titlepage
%    \end{macrocode}
% \end{macro}
%
% \begin{macro}{\whut@declare@mk}
% \begin{macro}{\whut@authorization@mk}
%    \begin{macrocode}
%%
%%
%% 独创性声明
\newcommand{\whut@declare@mk}{%
    \vspace*{-0.9cm}
    \renewcommand{\baselinestretch}{1.6}
    \begin{center}\zihao{-2}\heiti\ziju{0.4}{\whut@declaretitle}\end{center}\par
    \zihao{4}\ziju{0.1}{\whut@declare}\par
    \vskip0.8cm
    \begin{center}
        \whut@declaresigncap \underline{\makebox[3.5cm][s]{}}
        \whut@declaredatecap \underline{\makebox[3.5cm][s]{}}
    \end{center}
}
%% 论文授权
\newcommand{\whut@authorization@mk}{%
    \vskip1cm
    \renewcommand{\baselinestretch}{1.6}
    \begin{center}\zihao{-2}\heiti\whut@authorizationtitle\end{center}
    \zihao{4}\par\whut@authorization \par
    \begin{center} \whut@authorizationaddon \end{center}\par
    %%\vskip1cm
    \noindent\whut@authorsig\hspace*{2cm}\whut@teachersig\hspace*{2cm}\whut@frontdate
}
%    \end{macrocode}
% \end{macro}
% \end{macro}
%
%
% \begin{macro}{\makecover}
%    \begin{macrocode}
%% 生成封面
\newcommand{\makecover}{
    \phantomsection
    \pdfbookmark[-1]{\whut@ctitle}{ctitle}
    \normalsize
    \begin{titlepage}
% 论文封面和封二
        \whut@first@coverpage
        \clearpage
        \thispagestyle{empty}
        \whut@first@titlepage
% 独创性声明、授权说明
        \clearpage
        \thispagestyle{empty}
        \whut@declare@mk
        \whut@authorization@mk
    \end{titlepage}
}
%</cls>
%    \end{macrocode}
% \end{macro}
%
% \subsubsection{摘要}
% \label{sec:abstract}
% \begin{environment}{cabstract}
% \begin{environment}{eabstract}
%<*cls>
%% 为摘要的标题格式问题折腾了好久，几欲放弃，原来是没有正确设置好CTEXsetup的参数，
%% 在这里终于感觉到ctexbook的方便了。
%    \begin{macrocode}
\newenvironment{cabstract}{%
    \CTEXsetup[titleformat={\zihao{-3}\bfseries},beforeskip={0pt},afterskip={20pt}]{chapter}
    \chapter{\cabstractname}
    \zihao{-4}
    \@afterheading}
    {\par\vspace{2em}\par}
\newenvironment{eabstract}{%
    \CTEXsetup[titleformat={\zihao{-3}\bfseries},beforeskip={0pt},afterskip={20pt}]{chapter}
    \chapter{\eabstractname}
    \zihao{-4}
    \@afterheading}
    {\par\vspace{2em}\par}
%    \end{macrocode}
%</cls>
% \end{environment}
% \end{environment}
%
% \begin{macro}{\ckeywords}
% \begin{macro}{\ekeywords}
%<*cls>
%% 关键字
%    \begin{macrocode}
\newcommand\ckeywords[1]{{\heiti\zihao{-4} \whut@ckeywords@title} #1}
\newcommand\ekeywords[1]{\textbf{\zihao{-4} \whut@ekeywords@title} #1}
%    \end{macrocode}
%</cls>
% \end{macro}
% \end{macro}
%
% \subsubsection{主要符号表}
% \label{sec:denotationfmt}
% \begin{environment}{denotation}
%    \begin{macrocode}
%<*cls>
%% 主要符号对照表
\newenvironment{denotation}[1][2.5cm]{
    \chapter{\whut@denotation@name} % no tocline
    \noindent\begin{list}{}%
    {\vskip-30bp\zihao{-4}
     \renewcommand\makelabel[1]{##1\hfil}
     \setlength{\labelwidth}{#1} % 标签盒子宽度
     \setlength{\labelsep}{0.5cm} % 标签与列表文本距离
     \setlength{\itemindent}{0cm} % 标签缩进量
     \setlength{\leftmargin}{\labelwidth+\labelsep} % 左边界
     \setlength{\rightmargin}{0cm}
     \setlength{\parsep}{0cm} % 段落间距
     \setlength{\itemsep}{0cm} % 标签间距
    \setlength{\listparindent}{0cm} % 段落缩进量
    \setlength{\topsep}{0pt} % 标签与上文的间距
   }}{\end{list}}
%</cls>
%    \end{macrocode}
% \end{environment}
%
%
% \subsubsection{致谢以及声明}
% \label{sec:ackanddeclare}
%
% \begin{environment}{ack}
%    \begin{macrocode}
%<*cls>
%% 致谢
\newenvironment{ack}{%
%%    \clearpage
    \CTEXsetup[titleformat={\zihao{-2}\bfseries},beforeskip={0pt},afterskip={20pt}]{chapter}
    \chapter*{\whut@ackname}
    \addcontentsline{toc}{chapter}{\whut@ackname}
    \@afterheading
    \kaishu\zihao{-4}
}{\par\vspace{2em}\par}
%</cls>
%    \end{macrocode}
% \end{environment}
%
% \subsubsection{索引部分}
% \label{sec:threeindex}
% \begin{macro}{\listoffigures}
% \begin{macro}{\listoffigures*}
% \begin{macro}{\listoftables}
% \begin{macro}{\listoftables*}
%    \begin{macrocode}
%<*cls>
%% 索引
\def\whut@starttoc#1{% #1: float type, prepend type name in \listof*** entry.
    \let\oldnumberline\numberline
    \def\numberline##1{\oldnumberline{\csname #1name\endcsname\hskip.4em ##1}}
    \@starttoc{\csname ext@#1\endcsname}
    \let\numberline\oldnumberline}
\def\whut@listof#1{% #1: float type
    \@ifstar
    {\chapter[]{\csname list#1name\endcsname}\whut@starttoc{#1}}
    {\chapter{\csname list#1name\endcsname}\whut@starttoc{#1}}}
\renewcommand\listoffigures{\whut@listof{figure}}
\renewcommand*\l@figure{\@dottedtocline{1}{0em}{4em}}
\renewcommand\listoftables{\whut@listof{table}}
\let\l@table\l@figure
%    \end{macrocode}
% \end{macro}
% \end{macro}
% \end{macro}
% \end{macro}
%
% \begin{macro}{\equcaption}
%   本命令只是为了生成公式列表，所以这个 caption 是假的。如果要编号最好用
%    equation 环境，如果是其它编号环境，请手动添加添加 \cs{equcaption}。
% 用法如下：
%
% \cs{equcaption}\marg{counter}
%
% \marg{counter} 指定出现在索引中的编号，一般取 \cs{theequation}，如果你是用
%  \pkg{amsmath} 的 \cs{tag}，那么默认是 \cs{tag} 的参数；除此之外可能需要你
% 手工指定。
%
%    \begin{macrocode}
\def\ext@equation{loe}
\def\equcaption#1{%
\addcontentsline{\ext@equation}{equation}%
    {\protect\numberline{#1}}}
%    \end{macrocode}
% \end{macro}
%
% \begin{macro}{\listofequations}
% \begin{macro}{\listofequations*}
% \LaTeX{}默认没有公式索引，此处定义自己的 \cs{listofequations}。
%    \begin{macrocode}
\newcommand\listofequations{\whut@listof{equation}}
\let\l@equation\l@figure
%</cls>
%    \end{macrocode}
% \end{macro}
% \end{macro}
%
%
% \subsubsection{参考文献}
% \label{sec:ref}
%
%    \begin{macrocode}
%<*cls>
%% 重定义参考文献样式，来自thu
\renewenvironment{thebibliography}[1]{%
   \chapter*{\bibname}%
   \zihao{5}
   \list{\@biblabel{\@arabic\c@enumiv}}%
        {\renewcommand{\makelabel}[1]{##1\hfill}
         \settowidth\labelwidth{1.1cm}
         \setlength{\labelsep}{0.4em}
         \setlength{\itemindent}{0pt}
         \setlength{\leftmargin}{\labelwidth+\labelsep}
         \addtolength{\itemsep}{-0.7em}
         \usecounter{enumiv}%
         \let\p@enumiv\@empty
         \renewcommand\theenumiv{\@arabic\c@enumiv}}%
    \sloppy\frenchspacing
    \clubpenalty4000%
    \@clubpenalty \clubpenalty
    \widowpenalty4000%
    \interlinepenalty4000%
    \sfcode`\.\@m}
   {\def\@noitemerr
     {\@latex@warning{Empty `thebibliography' environment}}%
    \endlist\frenchspacing}
%</cls>
%    \end{macrocode}
%
%
% \begin{macro}{\onlinecite}
% 正文引用模式。依赖于 \pkg{natbib} 宏包，修改其中的命令。
%    \begin{macrocode}
%<*cls>
%% 引用格式
\bibpunct{[}{]}{,}{s}{}{,}
\renewcommand\NAT@citesuper[3]{\ifNAT@swa%
    \unskip\kern\p@\textsuperscript{\NAT@@open #1\NAT@@close}%
    \if*#3*\else\ (#3)\fi\else #1\fi\endgroup}
\DeclareRobustCommand\onlinecite{\@onlinecite}
\def\@onlinecite#1{\begingroup\let\@cite\NAT@citenum\citep{#1}\endgroup}
%</cls>
%    \end{macrocode}
% \end{macro}
%
%
% \subsubsection{附录}
% \label{sec:appendix}
%
% \subsubsection{个人简历}
%
% 定义个人简历章节标题
% \begin{environment}{resume}
%<*cls>
%% 个人简历发表文章等。
%    \begin{macrocode}
%% 简历
\newenvironment{resume}{%
  % \cleardoublepage
    \chapter*{\whut@resume@title}
    \addcontentsline{toc}{chapter}{\whut@resume@title}
    }{}
%    \end{macrocode}
%</cls>
% \end{environment}
%
% \begin{macro}{\resumeitem}
% 个人简历里面会出现的以发表文章，在投文章等。
%    \begin{macrocode}
%<*cls>
\newcommand{\resumeitem}[1]{\vspace{2.5em}{\zihao{4}\heiti\centerline{#1}}\par}
%</cls>
%    \end{macrocode}
% \end{macro}
%
% \subsubsection{索引}
%
%    \begin{macrocode}
%<*cls>
%% 生成索引的一些命令，虽然我们暂时还用不到。
\iffalse
\newcommand{\bs}{\symbol{'134}}%Print backslash
% \newcommand{\bs}{\ensuremath{\mathtt{\backslash}}}%Print backslash
% Index entry for a command (\cih for hidden command index
\newcommand{\cih}[1]{%
  \index{commands!#1@\texttt{\bs#1}}%
  \index{#1@\texttt{\hspace*{-1.2ex}\bs #1}}}
\newcommand{\ci}[1]{\cih{#1}\texttt{\bs#1}}
% Package
\newcommand{\pai}[1]{%
  \index{packages!#1@\textsf{#1}}%
  \index{#1@\textsf{#1}}%
  \textsf{#1}}
% Index entry for an environment
\newcommand{\ei}[1]{%
  \index{environments!\texttt{#1}}%
  \index{#1@\texttt{#1}}%
  \texttt{#1}}
% Indexentry for a word (Word inserted into the text)
\newcommand{\wi}[1]{\index{#1}#1}
\fi
%</cls>
%    \end{macrocode}
%
% \subsubsection{自定义命令和环境}
% \label{sec:userdefine}
%
% \begin{macro}{\pozhehao}
% 定义破折号。两个字宽，ex 差不多是当前字体的一半高度，所以通过 \cs{rule} 可以简单
% 的完成破折号绘制。
%    \begin{macrocode}
%<*cls>
\newcommand{\pozhehao}{\kern0.3ex\rule[0.8ex]{2em}{0.1ex}\kern0.3ex}
%</cls>
%    \end{macrocode}
% \end{macro}
%
%    \begin{macrocode}
%<*cls>
%% 在模板文档结束时即装入配置文件，这样用户就能在导言区进行相应的修改，否则必须在 document 开始后才能，感觉不好。
\AtEndOfClass{%%
%% This is file `WHUTthesis.cfg',
%% generated with the docstrip utility.
%%
%% The original source files were:
%%
%% WHUTthesis.dtx  (with options: `cfg')
%% This is a generated file.
%% 
%% Copyright (C) 2010 by Hu,Weiyi <huweiyi@whut.edu.cn>
%% 
%% This file may be distributed and/or modified under the
%% conditions of the LaTeX Project Public License, either version 1.3a
%% of this license or (at your option) any later version.
%% The latest version of this license is in:
%% 
%% http://www.latex-project.org/lppl.txt
%% 
%% and version 1.3a or later is part of all distributions of LaTeX
%% version 2004/10/01 or later.
%% 
%% This is the configuration file of the WHUTthesis package with LaTeX2e.

\ProvidesFile{WHUTthesis.cfg}
[2010/03/02 v0.1.0 Thesis Template for Wuhan University of Technology]
\theoremsymbol{\ensuremath{\square}}
\newtheorem*{proof}{证明}
\theoremstyle{plain}
\theoremsymbol{}
\theoremseparator{：}
\newtheorem{assumption}{假设}[chapter]
\newtheorem{definition}{定义}[chapter]
\newtheorem{proposition}{命题}[chapter]
\newtheorem{lemma}{引理}[chapter]
\newtheorem{theorem}{定理}[chapter]
\newtheorem{axiom}{公理}[chapter]
\newtheorem{corollary}{推论}[chapter]
\newtheorem{exercise}{练习}[chapter]
\newtheorem{example}{例}[chapter]
\newtheorem{remark}{注释}[chapter]
\newtheorem{problem}{问题}[chapter]
\newtheorem{conjecture}{猜想}[chapter]
\renewcommand\contentsname{目\hspace{1em}录}
\renewcommand\listfigurename{插图索引}
\renewcommand\listtablename{表格索引}
\newcommand\listequationname{公式索引}
\newcommand\equationname{公式}
\renewcommand\bibname{参考文献}
\renewcommand\indexname{索引}
\renewcommand\figurename{图}
\renewcommand\tablename{表}
\newcommand{\cabstractname}{摘\hspace{1em}要}
\newcommand{\eabstractname}{Abstract}
\let\CJK@todaysave=\today
\def\CJK@todaysmall@short{\the\year 年 \the\month 月}
\def\CJK@todaysmall{\CJK@todaysmall@short \the\day 日}
\CTEXdigits{\whut@yeartmp}{\the\year}
\CTEXnumber{\whut@monthtmp}{\the\month}
\def\CJK@todaybig@short{\whut@yeartmp \ 年 \whut@monthtmp \ 月}
\def\CJK@todaybig{\CJK@todaybig@short \CJKnumber{\the\day} 日}
\def\CJK@today{\CJK@todaysmall}
\renewcommand\today{\CJK@today}
\newcommand\CJKtoday[1][1]{%
  \ifcase#1\def\CJK@today{\CJK@todaysave}
    \or\def\CJK@today{\CJK@todaysmall}
    \or\def\CJK@today{\CJK@todaybig}
  \fi}
\cdate{\CJK@todaybig@short}
\edate{\ifcase \month \or January\or February\or March\or April\or May%
\or June\or July \or August\or September\or October\or November
\or December\fi\unskip,\ \ \the\year}
%% 封面等一些标题串
\def\whut@title@sep{：}
\def\whut@ckeywords@separator{，}
\def\whut@ekeywords@separator{,}
\def\whut@categorynumbercap{分类号}
\def\whut@secretlevelcap{密级} %% 注明论文密级为公开、内部、秘密或机密
\def\whut@UDCcap{UDC}
\def\whut@schoolcodecap{学校代码}
\def\whut@schoolcode{10497}
\def\whut@thesisname{学位论文}
\def\whut@ctitlecap{题目}
\def\whut@etitlecap{英文题目}
\def\whut@authornamecap{研究生姓名}
\def\whut@advisorcap{指导教师}
\def\whut@advisornamecap{姓名}
\def\whut@advisortitlecap{职称}
\def\whut@advisordegreecap{学位}
\def\whut@advisordepartmentcap{单位名称}
\def\whut@advisorzipcodecap{邮编}
\def\whut@coadvisorcap{副指导教师}
\def\whut@applycap{申请学位级别}
\def\whut@majorcap{学科专业名称}
\def\whut@submitdatecap{论文提交日期}
\def\whut@oraldefensedatecap{论文答辩日期}
\def\whut@authorizeorganizationcap{学位授予单位}
\def\whut@authorizedatecap{学位授予日期}
\def\whut@chairmancap{答辩委员会主席}
\def\whut@annotatorcap{评阅人}
\def\whut@title@sep{：}
\def\whut@secretlevel{秘密}
\def\whut@secretyear{\the\year}
\def\whut@schoolname{武汉理工大学}
\def\whut@apply{（申请\whut@cdegree 学位论文）}
\def\whut@cdepartmentcap{培养单位}
\def\whut@cmajortitle{学科专业}
\def\whut@cauthortitle{研究生}
\def\whut@cadvisortitle{指导教师}
\def\whut@masterheader{武汉理工大学硕士学位论文}
\def\whut@doctorheader{武汉理工大学博士学位论文}
\newcommand{\whut@declaretitle}{独创性声明}
\newcommand{\whut@declare}{%
本人声明，所呈交的论文是本人在导师指导下进行的研究工作及取得的研究成果。
尽我所知，除了文中特别加以标注和致谢的地方外，论文中不包含其他人已经发表
或撰写过的研究成果，也不包含为获得武汉理工大学或其他教育机构的学位或证书而使用过的材料。
与我一同工作的同志对本研究所做的任何贡献均已在论文中作了明确的说明并表示了谢意。
}
\def\whut@declaresigncap{签名：}
\def\whut@declaredatecap{日期：}
\newcommand{\whut@authorizationtitle}{学位论文使用授权书}
\newcommand{\whut@authorization}{%
本人完全了解武汉理工大学有关保留、使用学位论文的规定，
即学校有权保留并向国家有关部门或机构送交论文的复印件和电子版，
允许论文被查阅和借阅。
本人授权武汉理工大学可以将本学位论文的全部内容编入有关数据库进行检索，
可以采用影印、缩印或其他复制手段保存或汇编本学位论文。
同时授权经武汉理工大学认可的国家有关机构或论文数据库使用或收录本学位论文，
并向社会公众提供信息服务。
}
\newcommand{\whut@authorizationaddon}{（保密的论文在解密后应遵守此规定）}
\newcommand{\whut@authorsig}{研究生（签名）：}
\newcommand{\whut@teachersig}{导师（签名）：}
\newcommand{\whut@frontdate}{日期：}
\newcommand{\whut@ckeywords@title}{关键词：}
\newcommand{\whut@ekeywords@title}{Key words:}
\newcommand{\whut@denotation@name}{主要符号对照表}
\newcommand{\whut@ackshortname}{致谢}
\newcommand{\whut@ackname}{致\hspace{1em}谢}
\newcommand{\whut@resume@title}{个人简历、在学期间发表的学术论文与研究成果}
%% 封面、摘要、版权、致谢变量定义
\endinput
%%
%% End of file `WHUTthesis.cfg'.
}
%</cls>
%    \end{macrocode}
%
% \begin{macro}{\whut@setup@pdfinfo}
%    \begin{macrocode}
%<*cls>
%% 设置一些 pdf 文档信息，依赖于 \pkg{hyperref} 宏包。
\def\whut@setup@pdfinfo{%
    \hypersetup{%
        pdftitle={\whut@ctitle},
        pdfauthor={\whut@cauthor},
        pdfsubject={\whut@cdegree},
        pdfkeywords={\whut@ckeywords},
        pdfcreator={\whut@cauthor},
        pdfproducer={\whutthesis}}
}
%</cls>
%    \end{macrocode}
% \end{macro}
% \Finale
%
%
%
%
%
%
% \iffalse
%    \begin{macrocode}
%<*dtx-style>
\ProvidesPackage{dtx-style}

\RequirePackage{calc}
\RequirePackage{array,longtable}
\RequirePackage{fancybox,fancyvrb}
\RequirePackage{xcolor}

\RequirePackage{txfonts}
%\RequirePackage{CJKutf8}
\RequirePackage{ctex}
\RequirePackage{xltxtra}
\newif\ifdtx@UTF\dtx@UTFtrue % todo: review this piece of code, why it needs an
                             % \iffalse in the comment? where is the real mis-match?
\ifdtx@UTF\else
  \IfFileExists{CJKpunct.sty}%
               {\RequirePackage{CJKpunct}}%
               {\ClassWarning{whutthesis}{no CJKpunct.sty availiable!}}
\fi
%\RequirePackage{CJKspace}

%% 一些有趣的符号图案
\RequirePackage{marvosym}

\RequirePackage{hyperref}
\hypersetup{%
  unicode=false,
  CJKbookmarks=false,
  bookmarksnumbered=true,
  bookmarksopen=true,
  bookmarksopenlevel=1,
  breaklinks=true,
  linkcolor=blue,
  plainpages=false,
  pdfpagelabels,
  pdfborder=0 0 0}
\RequirePackage{url}
\RequirePackage{indentfirst}

% listings 源代码显示宏包
\RequirePackage{listings}
\lstset{tabsize=4, %
  frame=shadowbox, %把代码用带有阴影的框圈起来
  commentstyle=\color{red!50!green!50!blue!50},%浅灰色的注释
  rulesepcolor=\color{red!20!green!20!blue!20},%代码块边框为淡青色
  keywordstyle=\color{blue!90}\bfseries, %代码关键字的颜色为蓝色，粗体
  showstringspaces=false,%不显示代码字符串中间的空格标记
  stringstyle=\ttfamily, % 代码字符串的特殊格式
  keepspaces=true, %
  breakindent=22pt, %
  numbers=left,%左侧显示行号
  stepnumber=1,%
  numberstyle=\tiny, %行号字体用小号
  basicstyle=\footnotesize, %
  showspaces=false, %
  flexiblecolumns=true, %
  breaklines=true, %对过长的代码自动换行
  breakautoindent=true,%
  breakindent=4em, %
  escapebegin=\begin{CJK*}{GBK}{kai},escapeend=\end{CJK*},
  aboveskip=1em, %代码块边框
  %% added by http://bbs.ctex.org/viewthread.php?tid=53451
  fontadjust,
  captionpos=t,
  framextopmargin=2pt,framexbottommargin=2pt,abovecaptionskip=-3pt,belowcaptionskip=3pt,
  xleftmargin=4em,xrightmargin=4em, % 设定listing左右的空白
  texcl=true,
  % 设定中文冲突，断行，列模式，数学环境输入，listing数字的样式
  extendedchars=false,columns=flexible,mathescape=true
  % numbersep=-1em
}
\renewcommand{\lstlistingname}{代码} %% 重命名Listings标题头

\setlength{\parskip}{4pt plus1pt minus0pt}
\setlength{\topsep}{0pt}
\setlength{\partopsep}{0pt}
\setlength{\parindent}{20pt}
\addtolength{\oddsidemargin}{-1cm}
\advance\textwidth 1.5cm
\addtolength{\topmargin}{-1cm}
\addtolength{\headsep}{0.3cm}
\addtolength{\textheight}{2.3cm}

\newcommand\hei{\CJKfamily{hei}}
\newcommand\fs{\CJKfamily{fs}}

\renewcommand{\baselinestretch}{1.3}
\setlength{\shadowsize}{3pt}
\def\DescribeOption#1{\SpecialOptionIndex{#1}}
\def\SpecialOptionIndex#1{\index{#1\actualchar\textbf{#1}}}
\renewenvironment{description}
  {\list{}{\setlength\labelwidth{2cm}%
           \setlength\labelsep{3pt}%
           \setlength\leftmargin{\labelwidth+\labelsep}%
           \addtolength{\itemsep}{3pt}%
           \renewcommand\makelabel[1]{%
             \shadowbox{\color{blue!90}\texttt##1}\DescribeOption{##1}}}
  }{\endlist}
\DefineVerbatimEnvironment{example}{Verbatim}%
  {frame=single,framerule=0.3mm,rulecolor=\color{red!75!green!50!blue},%
   fillcolor=\color{red!75!green!50!blue!15},framesep=2mm,baselinestretch=1.2,%
   fontsize=\small,gobble=1}
\DefineVerbatimEnvironment{shell}{Verbatim}%
  {frame=single,framerule=0.3mm,rulecolor=\color{red!85!green!60},%
   fillcolor=\color{red!85!green!10},framesep=2mm,fontsize=\small,gobble=1}
\long\def\myentry#1{\vskip5pt\par\noindent\llap{{\color{blue}\fs #1}}\marginpar{\strut}\hskip\parindent}
\def\tableofcontents{\renewcommand{\baselinestretch}{1.0}\@starttoc{toc}}
\def\DescribeMacro{\Describe@Macro}
\def\Describe@Macro#1{\PrintDescribeMacro{#1}\SpecialUsageIndex{#1}}
\def\PrintDescribeMacro#1{{\color{-red!75!green!50!blue!55}\MacroFont \string #1\hskip1em}}
\def\ps@headings{%
  \let\@oddfoot\@empty
  \def\@oddhead{\vbox{%
    \hb@xt@ \textwidth{\llap{\fbox{\rightmark\rule[-2pt]{0pt}{13pt}}}\hfil\thepage}%
    \vskip-0.7pt%
    \hb@xt@ \textwidth{\hrulefill}}}
  \let\@evenfoot\@oddfoot
  \let\@evenhead\@oddhead
  \let\@mkboth\markboth
  \def\sectionmark##1{%
    \markright{\ifnum \c@secnumdepth >\m@ne
      \thesection\quad
      \fi
      ##1}}
  \def\subsectionmark##1{%
    \markright{\ifnum \c@secnumdepth >\m@ne
      \thesubsection\quad
      \fi
      ##1}}
  \def\subsubsectionmark##1{%
    \markright{\ifnum \c@secnumdepth >\m@ne
      \thesubsubsection\quad
      \fi
      ##1}}}
\renewcommand\section{\@startsection{section}{1}{\z@}%
                                    {-3.5ex \@plus -1ex \@minus -.2ex}%
                                    {2.3ex \@plus.2ex}%
                                    {\normalfont\Large\bfseries\hei}}

\renewcommand\subsection{\@startsection{subsection}{2}{\z@}%
                                       {-3.25ex\@plus -1ex \@minus -.2ex}%
                                       {1.5ex \@plus .2ex}%
                                       {\normalfont\large\bfseries\hei}}
\renewcommand\subsubsection{\@startsection{subsubsection}{3}{\z@}%
                                          {-3.25ex\@plus -1ex \@minus -.2ex}%
                                          {1.5ex \@plus .2ex}%
                                          {\normalfont\normalsize\bfseries\hei}}
\renewcommand\paragraph{\@startsection{paragraph}{4}{\z@}%
                                      {3.25ex \@plus1ex \@minus.2ex}%
                                      {-1em}%
                                      {\normalfont\normalsize\bfseries}}
\renewcommand\subparagraph{\@startsection{subparagraph}{5}{\parindent}%
                                         {3.25ex \@plus1ex \@minus .2ex}%
                                         {-1em}%
                                         {\normalfont\normalsize\bfseries}}
\pagestyle{empty}
%</dtx-style>
%    \end{macrocode}
% \fi
%
\endinput
